\exercisesection{\S~3}

\begin{enumerate}
\def\labelenumi{\arabic{enumi})}
\item
  设 G 是维数 \textgreater{} 0 的紧李群。则 G 的每个有限交换子群都是循环群，当且仅当 G 同构于 U、SU(2, C) 或 SU(2, C) 中某个极大环面的正规化子。
\item
  用 K 表示域 R、C 或 H 之一，用 n 表示一个整数 \(\geq 1\)。给空间 \(K_{s}^{n}\) 赋以通常的埃尔米特型。
\end{enumerate}

\begin{enumerate}
\def\labelenumi{\alph{enumi})}
\item
  证明 \(\mathbf{U}(n,\mathrm{K})\) 是紧实李群。
\item
  证明 \(\mathrm{K}_s^n\) 中半径为 1 的球面是 \(\mathbf{U}(n,\mathbf{K})\) 的（实）李齐性空间；一点的稳定子群同构于 \(\mathbf{U}(n - 1,\mathbf{K})\)。
\item
  推出群 \(\mathbf{U}(n,\mathbf{C})\) 与 \(\mathbf{U}(n,\mathbf{H})\) 都连通，而 \(\mathbf{O}(n,\mathbf{R})\) 有两个连通分支。
\item
  证明群 \(\mathbf{SO}(n,\mathbf{R})\) 与 \(\mathbf{SU}(n,\mathbf{C})\) 都连通。
\item
  证明群 \(\mathbf{O}(n,\mathbf{R})\)（相应地 \(\mathbf{U}(n,\mathbf{C})\)）是 Z/2Z（相应地 T）借助于 \(\mathbf{SO}(n,\mathbf{R})\)（相应地 \(\mathbf{SU}(n,\mathbf{C})\)）的半直积。
\end{enumerate}

\begin{enumerate}
\def\labelenumi{\arabic{enumi})}
\setcounter{enumi}{2}
\tightlist
\item
  \begin{enumerate}
  \def\labelenumii{\alph{enumii})}
  \tightlist
  \item
    证明实李群 \(\mathbf{U}(n, \mathbf{H})\) 的李代数是由诸如 \(x \in \mathrm{M}_n(\mathbf{H})\) 且满足 \(^t\bar{x} = -x\) 的矩阵组成的集合，括号为 \([x, y] = xy - yx\)。把它记作 \(\mathfrak{u}(n, \mathbf{H})\)。
  \end{enumerate}
\end{enumerate}

\begin{enumerate}
\def\labelenumi{\alph{enumi})}
\setcounter{enumi}{1}
\item
  把 \(\mathbf{C}\) 与子域 \(\mathbf{R}(i)\) 等同起来（它是 \(\mathbf{H}\) 的一个子域），并把 \(\mathbf{C}^{2n}\) 与 \(\mathbf{H}^n\) 借助同构 \((z_1, \ldots, z_{2n}) \mapsto (z_1 + jz_{n+1}, \ldots, z_n + jz_{2n})\) 等同起来。证明等式 \(\mathbf{U}(n, \mathbf{H}) = \mathbf{U}(2n, \mathbf{C}) \cap \mathbf{Sp}(2n, \mathbf{C})\)。
\item
  由 \(b)\) 推出，每个 \(\mathbf{C}_n\) 型的紧单（实）李代数都同构于 \(\mathfrak{u}(n, \mathbf{H})\)。
\end{enumerate}

\begin{enumerate}
\def\labelenumi{\arabic{enumi})}
\setcounter{enumi}{3}
\tightlist
\item
  设 \(n\) 是一个整数 \(\geq 1\)。
\end{enumerate}

\begin{enumerate}
\def\labelenumi{\alph{enumi})}
\item
  证明群 \(\mathbf{SU}(n,\mathbf{C})\) 是单连通的（利用习题 2）。
\item
  证明 \(\mathbf{SU}(n,\mathbf{C})\) 的中心由诸矩阵 \(\lambda .I_n\)（其中 \(\lambda \in \mathbf{C}\)，\(\lambda^n = 1\)）组成。
\item
  每个 \(A_{n}\) 型的概单紧李群都同构于 \(\mathbf{SU}(n+1,\mathbf{C})\) 关于由矩阵 \(\zeta^k. I_{n+1} (0 \leq k < d)\) 组成的循环子群的商群，其中 \(d\) 整除 \(n + 1\)，\(\zeta\) 是一个本原 \(d\) 次单位根。
\item
  证明群 \(\mathbf{SL}(n,\mathbf{C})\) 是单连通的（利用第 III 章，§6，节 9，定理 6）。
\end{enumerate}

\begin{enumerate}
\def\labelenumi{\arabic{enumi})}
\setcounter{enumi}{4}
\tightlist
\item
  设 \(n\) 为整数 \(\geq 1\)，用 \(\mathbf{Spin}(n, \mathbf{R})\) 表示对应于 \(\mathbf{R}^n\) 上通常二次型的既约 Clifford 群（《代数》第 IX 章，§9，节 5）。
\end{enumerate}

\begin{enumerate}
\def\labelenumi{\alph{enumi})}
\item
  证明 \(\mathbf{Spin}(n,\mathbf{R})\) 是紧李群，且满同态 \(\varphi:\mathbf{Spin}(n,\mathbf{R})\to\mathbf{SO}(n,\mathbf{R})\)（同上）是解析的，其核为 \(\{+1,-1\}\)。
\item
  证明当 \(n \geq 2\) 时，\(\mathbf{Spin}(n, \mathbf{R})\) 是连通且单连通的（利用习题 2）。群 \(\pi_{1}(\mathbf{SO}(n, \mathbf{R}))\) 是 2 阶循环群。
\item
  设 \(Z_{n}\) 是 \(\mathbf{Spin}(n,\mathbf{R})\) 的中心。证明当 n 为奇数时 \(Z_{n}=\{+1,-1\}\)，当 n 为偶数时 \(Z_{n}=\{+1,-1,\varepsilon,-\varepsilon\}\)，其中 \(\varepsilon=e_{1}\ldots e_{n}\) 是 \(R^{n}\) 的典则基中诸元素之积；当 r 为偶数（相应地奇数）时，群 \(Z_{2r}\) 同构于 \((\mathbf{Z}/2\mathbf{Z})^{2}\)（相应地 \(Z/4Z\)）。
\item
  证明每个 \(B_{n}\) 型（\(n \geq 2\)）的概单连通紧李群都同构于 \(\mathbf{Spin}(2n + 1, \mathbf{R})\) 或 \(\mathbf{SO}(2n + 1, \mathbf{R})\)。
\item
  若 r 为奇数（相应地偶数）且 \(\geq 2\)，则每个 \(D_{r}\) 型的概单连通紧李群同构于 \(\mathbf{Spin}(2r,\mathbf{R})\)、\(\mathbf{SO}(2r,\mathbf{R})\) 或 \(\mathbf{SO}(2r,\mathbf{R})/\{\pm I_{2r}\}\)（相应地同构于上述群之一或 \(\mathbf{Spin}(2r,\mathbf{R})/\{1,\varepsilon\}\)）。
\end{enumerate}

\begin{enumerate}
\def\labelenumi{\arabic{enumi})}
\setcounter{enumi}{5}
\tightlist
\item
  \begin{enumerate}
  \def\labelenumii{\alph{enumii})}
  \tightlist
  \item
    证明紧李群 \(\mathbf{U}(n,\mathbf{H})\) 连通且单连通（利用习题 2），且其中心为 \(\{\pm I_{n}\}\)。
  \end{enumerate}
\end{enumerate}

\begin{enumerate}
\def\labelenumi{\alph{enumi})}
\setcounter{enumi}{1}
\tightlist
\item
  每个 \(C_{n}\) 型（\(n \geq 3\)）的概单连通紧李群都同构于 \(\mathbf{U}(n, \mathbf{H})\) 或 \(\mathbf{U}(n, \mathbf{H}) / \{\pm I_{n}\}\)。
\end{enumerate}

\begin{enumerate}
\def\labelenumi{\arabic{enumi})}
\setcounter{enumi}{6}
\tightlist
\item
  设 A 是 Cayley 八元数代数（《代数》第 III 章，附录，节 3），并取同上中的基 \((e_i)_{0\leq i\leq 7}\)。用 V 表示由 \(e_1,\ldots ,e_7\) 生成的纯八元数子空间，用 E 表示 V 中由 \(e_1,e_2,e_3,e_5,e_6,e_7\) 生成的子空间。把 A 中由 \(e_0,e_4\) 生成的子代数与复数域 \(\mathbf{C}\) 等同起来，并用 G 表示幺元代数 A 的自同构拓扑群。
\end{enumerate}

\begin{enumerate}
\def\labelenumi{\alph{enumi})}
\item
  用 Q 表示由 Cayley 范数诱导的 V 上的二次型，使 \((e_{i})_{1\leq i\leq7}\) 是 V 的一个标准正交基。证明映射 \(\sigma\mapsto\sigma|V\) 是从 G 到群 \(\mathbf{SO}(\mathbf{Q})\) 的单同态，而该群同构于 \(\mathbf{SO}(7,\mathbf{R})\)。
\item
  证明 A 的乘法赋予 E 一个 3 维 C-向量空间结构，且 \(\{e_{1}, e_{2}, e_{3}\}\) 是它的一个基。用 \(\Phi\) 表示 E 上以该基为标准正交基的埃尔米特型。设 H 是 G 中 \(e_{4}\) 的稳定子群。证明映射 \(\sigma \mapsto \sigma|E\) 是从 H 到群 \(\mathbf{SU}(\Phi)\) 的同构，而该群同构于 \(\mathbf{SU}(3, \mathbf{C})\)。映射 \(\sigma \mapsto \sigma(e_{4})\) 诱导出 G/H 到 V 的球面中的一个嵌入，该球面同构于 \(S_{6}\)。
\item
  设 T 是 H 中由自同构 \(\sigma\)（满足 \(\sigma(e_{0}) = e_{0}\)，\(\sigma(e_{1}) = \alpha e_{1}\)，\(\sigma(e_{2}) = \beta e_{2}\)，\(\sigma(e_{3}) = \gamma e_{3}\)，\(\sigma(e_{4}) = e_{4}\)，\(\sigma(e_{5}) = \bar{\alpha} e_{5}\)，\(\sigma(e_{6}) = \bar{\beta} e_{6}\)，\(\sigma(e_{7}) = \bar{\gamma} e_{7}\)）组成的环面，其中 \(\alpha, \beta, \gamma\) 是三个模为 1 且满足 \(\alpha\beta\gamma = 1\) 的复数。设 N 是 T 在 G 中的正规化子；证明 N/T 的阶为 12（注意 N 的每个元素都必须使诸 \(\pm e_{i}, i \neq 0\) 组成的集合稳定）；由此推出 G 的秩为 2。
\item
  证明 G 是 \(G_{2}\) 型的连通半单群，且 G/H 可与 \(S_{6}\) 等同（证明 \(G_{0} \neq H\)；推出 \(G_{0}\) 是 \(G_{2}\) 型的，再用维数论证）。每个 \(G_{2}\) 型紧群都同构于 G。
\end{enumerate}

\begin{enumerate}
\def\labelenumi{\arabic{enumi})}
\setcounter{enumi}{7}
\item
  设 G 是 \(A_{n}, B_{n}, C_{n}, D_{n}\) 或 \(G_{2}\) 型的概单连通紧李群。证明 \(\pi_{2}(G) = 0\) 且 \(\pi_{3}(G) = Z\)（利用习题 2 至 7，以及 \(\pi_{i}(\mathbf{S}_{n})\) 当 i \textless{} n 时为零、当 i = n 时为循环群这一事实（参见~《一般拓扑学》第 XI 章）。
\item
  设 g 是紧李代数，t 是 g 的一个嘉当子代数；设 \((\mathrm{X}_{\alpha})_{\alpha\in\mathbb{R}}\) 是分裂约化代数 \((\mathfrak{g}_{\mathbf{C}},\mathfrak{t}_{\mathbf{C}})\) 中的一个谢瓦莱系，使得 \(X_{\alpha}\) 与 \(X_{-\alpha}\)（关于 g）对所有 \(\alpha\in R\) 共轭。
\end{enumerate}

\begin{enumerate}
\def\labelenumi{\alph{enumi})}
\item
  设 \(\mathcal{T}\) 是一个 \(\mathbf{Z}\)-子模，包含于 \(\mathfrak{t}_{\mathbf{C}}\)，它包含诸 \(iH_{\alpha}\)（\(\alpha \in \mathrm{R}\)），且满足 \(\alpha(\mathcal{T}) \subset \mathbf{Z}i\)（对所有 \(\alpha \in \mathrm{R}\)）。证明 \(\mathbf{Z}\)-子模 \(\mathcal{G}\)（\(\mathfrak{g}_{\mathbf{C}}\) 中由 \(\mathcal{T}\) 和诸元素 \(u_{\alpha}\)、\(v_{\alpha}\)（\(\alpha \in \mathrm{R}\)）生成的）是一个 \(\mathbf{Z}\)-李子代数（在 \(\mathfrak{g}\) 中）。
\item
  设 g（相应地 t）是紧群 G（相应地 G 的一个极大环面 T）的李代数。设 \(\Gamma(\mathrm{T})\) 是同态 \(\exp_{\mathrm{T}}:t\to\mathrm{T}\) 的核。证明 Z-模 \((2\pi)^{-1}\varGamma(\mathrm{T})\) 满足 a) 的假设。
\item
  设 \(\langle,\rangle\) 是 g 上的一个不变内积；设 \(\mu\)（相应地 \(\tau\)） 是当 g（相应地 t）借助一个标准正交基与空间 \(R^{n}\) 等同时，对应于 Lebesgue 测度的 g（相应地 t）上的 Haar 测度。再用 \(\mu\)（相应地 \(\tau\)） 表示 g/G（相应地 t/T）上由 \(\mu\)（相应地 \(\tau\)） 除以 G（相应地 T）上的规范化 Haar 测度所得的商测度。证明公式 \(\mu(\mathfrak{g}/\mathcal{G})=\tau(\mathfrak{t}/\mathcal{T}).\prod_{\alpha\in\mathrm{R}_{+}}\langle iH_{\alpha},iH_{\alpha}\rangle\)。
\end{enumerate}

\exercisesection{\S~4}

\begin{enumerate}
\def\labelenumi{\arabic{enumi})}
\tightlist
\item
  取 G 为群 \(\mathbf{SU}(n,\mathbf{C})\) 或 \(\mathbf{U}(n,\mathbf{H})\) 之一；把 \(g_{C}\) 相应地与 \(\mathfrak{sl}(n,\mathbf{C})\) 或 \(\mathfrak{sp}(2n,\mathbf{C})\) 等同，参见~§3，节 4 及习题 3。我们采用第 VIII 章，§13，节 1 与节 3 的记号，其中 k=C。
\end{enumerate}

\begin{enumerate}
\def\labelenumi{\alph{enumi})}
\item
  证明 G 中由复系数对角矩阵组成的子群 T 是一个极大环面，且 \(\mathrm{L}(\mathrm{T})_{(\mathbf{C})} = \mathfrak{h}\)。
\item
  把 X(T) 与 \(\mathfrak{h}^*\) 的一个子群等同起来（借助同态 \(\delta\)）。证明线性形式 \(\varepsilon_i\) 与支配权 \(\varpi_j\) 都属于 X(T)。若 \(t = \mathrm{diag}(t_1, \ldots, t_n) \in \mathrm{T}\)，则 \(\varepsilon_i(t) = t_i\) 且 \(\varpi_j(t) = t_1 \ldots t_j\)（对 \(1 \leq i, j \leq n\)）。
\item
  由 b) 给出群 \(\mathbf{SU}(n,\mathbf{C})\) 与 \(\mathbf{U}(n,\mathbf{H})\) 单连通（参见~§3，习题 4 与 6）这一事实的另一证明。
\end{enumerate}

\begin{enumerate}
\def\labelenumi{\arabic{enumi})}
\setcounter{enumi}{1}
\tightlist
\item
  取 \(\mathrm{G} = \mathbf{SO}(n, \mathbf{R})\)，其中 \(n \geq 3\)；当 n 为奇数时置 \(n = 2l + 1\)，当 n 为偶数时置 n = 2l。代数 \(g_{C}\) 可与 \(\mathfrak{o}(n, \mathbf{C})\) 等同；我们采用第 VIII 章，§13，节 2 与节 4 的记号。用 \((f_i)_{1 \leq i \leq n}\) 表示 \(R^n\) 的典则基。置 \(e_j = \frac{1}{\sqrt{2}}(f_{2j-1} + if_{2j})\) 与 \(e_{-j} = \frac{1}{\sqrt{2}}(f_{2j-1} - if_{2j})\)（对 \(1 \leq j \leq l\)），当 n 为奇数时再置 \(e_0 = i\sqrt{2}f_{2l+1}\)；当 n 为偶数（相应地奇数）时，选取 \(C^n\) 的由 \(e_1, \ldots, e_l, e_{-l}, \ldots, e_{-1}\)（相应地 \(e_1, \ldots, e_l, e_0, e_{-l}, \ldots, e_{-1}\)）给出的 Witt 基。
\end{enumerate}

\begin{enumerate}
\def\labelenumi{\alph{enumi})}
\tightlist
\item
  设 \(H_{i}\) 是 \(R^{n}\) 中由 \(f_{2i-1}\) 与 \(f_{2i}\) 生成的子空间（\(1 \leq i \leq l\)）；证明 G 中由满足 \(g(\mathrm{H}_{i}) \subset \mathrm{H}_{i}\) 且 \(\det(g|\mathrm{H}_{i}) = 1\)（对 \(1 \leq i \leq l\)）的元素 g 组成的子群是 G 的一个极大环面 T，且 \(\mathrm{L}(\mathrm{T})_{(\mathbf{C})} = \mathfrak{h}\)。
\item
  把 X(T) 与 \(h^{*}\) 的一个子群等同起来（借助 \(\delta\)）。证明线性形式 \(\varepsilon_{i}\) 属于 X(T)；若 \(t \in T\) 且 t 在 \(H_{j}\) 上的限制是转角为 \(\theta_{j}\) 的旋转，则 \(\varepsilon_{j}(t) = e^{i\theta_{j}}\)。诸权 \(\varpi_{1}, \ldots, \varpi_{l-2}\)；\(2\varpi_{l-1}, 2\varpi_{l}, \varpi_{l-1} \pm \varpi_{l}\) 都属于 X(T)。若 n 为奇数，则 \(\varpi_{l-1}\) 属于 X(T)。
\item
  设 \(\tilde{\mathbf{G}} = \mathbf{Spin}(n, \mathbf{R})\)，并设 \(\varphi : \tilde{G} \to G\) 是典则覆盖。对 \(\boldsymbol{\theta} = (\theta_{1}, \ldots, \theta_{l}) \in \mathbf{R}^{l}\)，置 \(t(\boldsymbol{\theta}) = \prod_{i=1}^{l} (\cos\theta_{i} - f_{2i-1}f_{2i}\sin\theta_{i}) \in \tilde{\mathbf{G}}\)。证明诸 \(t(\boldsymbol{\theta})\)（\(\theta \in R^{l}\)）组成的集合是一个极大环面 \(\tilde{T}\)（\(\tilde{G}\) 中的），且 \(\varphi(\tilde{T}) = T\)。若把 X( \(\tilde{T}\) ) 与 \(h^{*}\) 的一个子群等同起来，则有 \(\varepsilon_{j}(t(\boldsymbol{\theta})) = e^{2i\theta_{j}}\)。
\item
  证明权 \(\varpi_{l-1}\) 与 \(\varpi_{l}\) 属于 X( \(\tilde{T}\) )；由此推出 \(\mathbf{Spin}(n,\mathbf{R})\) 是单连通的（参见~§3，习题 5）。
\end{enumerate}

\begin{enumerate}
\def\labelenumi{\arabic{enumi})}
\setcounter{enumi}{2}
\tightlist
\item
  \begin{enumerate}
  \def\labelenumii{\alph{enumii})}
  \tightlist
  \item
    证明 \(\sigma: \mathrm{A} \mapsto \overline{\mathrm{A}}\) 这一 \(\mathbf{SU}(n, \mathbf{C})\) 的自同构当 \(n \geq 3\) 时不是内自同构。\(\mathbf{SU}(n, \mathbf{C})\) 的每个非内自同构都具有 \((\operatorname{Int} g) \circ \sigma\)（\(g \in \mathbf{SU}(n, \mathbf{C})\)）的形式。
  \end{enumerate}
\end{enumerate}

\begin{enumerate}
\def\labelenumi{\alph{enumi})}
\setcounter{enumi}{1}
\tightlist
\item
  证明对每个 \(A_{n}\) 型（\(n \geq 2\)）的概单连通紧群 G，群 \(\operatorname{Aut}(\mathrm{G})/\operatorname{Int}(\mathrm{G})\) 都是 2 阶循环群（参见~§3，习题 4）。
\end{enumerate}

\begin{enumerate}
\def\labelenumi{\arabic{enumi})}
\setcounter{enumi}{3}
\tightlist
\item
  设 \(n\) 为整数 \(\geq 2\)。
\end{enumerate}

\begin{enumerate}
\def\labelenumi{\alph{enumi})}
\item
  设 \(g \in \mathbf{O}(2n, \mathbf{R})\) 且 \(\det g = -1\)；证明 \(\operatorname{Int} g\) 这一 \(\mathbf{O}(2n, \mathbf{R})\) 的自同构诱导出 \(\mathbf{SO}(2n, \mathbf{R})\) 的一个非内自同构。
\item
  当 \(n \geq 2\) 时，群 \(\operatorname{Aut}(\mathbf{SO}(2n, \mathbf{R}))\) 等于 \(\operatorname{Int}(\mathbf{O}(2n, \mathbf{R}))\)（它同构于 \(\mathbf{O}(2n, \mathbf{R}) / \{\pm I_{2n}\}\)）。
\item
  对 \(\mathbf{Spin}(2n,\mathbf{R})\) 与 \(\mathbf{SO}(2n,\mathbf{R})/\{\pm I_{2n}\}\) 建立类似的结果。若 n 为偶数且 \(\neq 2\)，则群 \(\mathbf{Spin}(2n,\mathbf{R})/\{1,\varepsilon\}\)（§3，习题 5）的每个自同构都是内自同构。
\end{enumerate}

\begin{enumerate}
\def\labelenumi{\arabic{enumi})}
\setcounter{enumi}{4}
\tightlist
\item
  设 \(\mathbf{R}\) 是一个既约且不可约的根系。
\end{enumerate}

\begin{enumerate}
\def\labelenumi{\alph{enumi})}
\item
  证明 R 中长度最大的根组成的集合是 R 的一个对称闭子集，且在 W(R) 下稳定。
\item
  证明 R 的每个非空、闭、对称且在 W(R) 下稳定的子集 P 都等于 R，或等于 R 中长度最大的根组成的集合。
\item
  设 \(P \neq R\)。证明根系 P 为 \(D_{l}\) 型（相应地 \((\mathrm{A}_{1})^{l}, \mathrm{D}_{4}, \mathrm{A}_{2}\)）当 R 为 \(B_{l}\) 型（相应地 \(C_{l}, F_{4}, G_{2}\)）时。
\end{enumerate}

\begin{enumerate}
\def\labelenumi{\arabic{enumi})}
\setcounter{enumi}{5}
\tightlist
\item
  设 \(\mathrm{H}\) 是 \(\mathrm{G}\) 的包含 \(\mathrm{N}_{\mathrm{G}}(\mathrm{T})\) 的闭子群。
\end{enumerate}

\begin{enumerate}
\def\labelenumi{\alph{enumi})}
\item
  证明 H 等于它在 G 中的正规化子，且 \(H/H_{0}\) 同构于 \(W/W_{H_{0}}(T)\)。
\item
  证明 \(\mathrm{R}(\mathrm{H}_{0},\mathrm{T})\) 在 W 下稳定。反之，若 K 是包含 T 的连通闭子群，且 \(\mathrm{R}(\mathrm{K},\mathrm{T})\) 在 W 下稳定，则 K 的正规化子包含 T 的正规化子。
\end{enumerate}

\(c)\) 设 \(\mathbf{G}\) 是概单的。证明 \(\mathbf{H}\) 等于 \(\mathrm{N_G(T)}\)，或等于 \(\mathbf{G}\)，或者（在同构意义下）处于下列情形之一：

\(\alpha\)（相应地 \(\alpha'\)）\(\mathrm{G} = \mathbf{Spin}(2l + 1, \mathbf{R})\)（相应地 \(\mathrm{G} = \mathbf{SO}(2l + 1, \mathbf{R})\)），且 \(\mathrm{H}_0\) 是 \(\mathbf{R}^{2l + 1}\) 中某个非零向量的稳定子群，它同构于 \(\mathbf{Spin}(2l, \mathbf{R})\)（相应地 \(\mathbf{SO}(2l, \mathbf{R})\)）；

\(\beta)\) \(\mathrm{G} = \mathrm{U}(l,\mathbf{H})\)，且 \(\mathrm{H}_0\) 是由对角矩阵组成的子群 D；\(\beta^{\prime})\) \(\mathrm{G} = \mathrm{U}(l,\mathbf{H}) / \{\pm 1\}\)，且 \(\mathrm{H}_0 = \mathrm{D} / \{\pm 1\}\)；

\(\gamma)\) G 是 \(\mathrm{F}_4\) 型的，且 \(\mathrm{H}_0\) 同构于 Spin(8,R)；

\(\delta)\) G 是 \(\mathrm{G}_2\) 型的，且 \(\mathrm{H}_0\) 是 §3 习题 7 中定义的子群（它同构于 \(\mathbf{SU}(3,\mathbf{C})\)）。

\begin{enumerate}
\def\labelenumi{\arabic{enumi})}
\setcounter{enumi}{6}
\tightlist
\item
  设 \(\tau : \mathbf{G} \to \mathbf{GL}(\mathbf{V})\) 是 \(\mathbf{G}\) 在一个有限维实向量空间上的连续表示。假定对所有 \(\lambda \in \mathrm{X}(\mathrm{T})\)，都有 \(\dim_{\mathbf{C}} \hat{V}_{\lambda} \leq 1\)。
\end{enumerate}

\begin{enumerate}
\def\labelenumi{\alph{enumi})}
\item
  证明表示 \(\tau\) 是一个由两两不同构的不可约表示组成的有限族 \((\tau_{i})_{1\leq i\leq s}\) 的直和，且其交换子 \(K_{i}\ (1\leq i\leq s)\) 同构于 R 或 C。
\item
  在 V 上存在使 \(\tau(\mathrm{G})\) 的运算均为 C-线性的 C-向量空间结构，当且仅当 \(K_{1},\ldots,K_{s}\) 都同构于 C；此时这种结构共有 \(2^{s}\) 个。
\end{enumerate}

\begin{enumerate}
\def\labelenumi{\arabic{enumi})}
\setcounter{enumi}{7}
\tightlist
\item
  设 H 是 G 的一个具有极大秩的连通闭子群，h 是它的李代数，X 是流形 G/H，V 是 X 在对应于 H 的类的点处的切空间；把 V 与商 g/h 等同起来。
\end{enumerate}

\begin{enumerate}
\def\labelenumi{\alph{enumi})}
\tightlist
\item
  若 j 是 X 上的一个殆复结构（《微分流形与解析流形，结果》，8.8.3），用 \(V''(j)\) 表示复向量空间 \(C \otimes V\) 中由满足 \(j(u) = -iu\) 的元素 u 组成的子空间，并用 \(\mathfrak{q}(j)\) 表示 \(g_{C}\) 中作为 \(V''(j)\) 的逆像的子空间，使得当 V 赋以由 j 诱导的 C-向量空间结构时，典则映射 \(V \to \mathfrak{g}_{\mathbf{C}} / \mathfrak{q}(j)\) 是一个 C-线性同构。证明映射 \(j \mapsto \mathfrak{q}(j)\) 是从 X 上在 G 下不变的殆复结构的集合，到 \(g_{C}\) 中满足下列条件的复子空间 p 的集合的一个双射：
\end{enumerate}

\[
\mathfrak {p} + \bar {\mathfrak {p}} = \mathfrak {g} _ {\mathbf {C}}\tag{1}
\]

\[
\mathfrak {p} \cap \bar {\mathfrak {p}} = \mathfrak {h} _ {\mathbf {C}}\tag{2}
\]

\[
[ \mathfrak {h}, \mathfrak {p} ] \subset \mathfrak {p}.\tag{3}
\]

\begin{enumerate}
\def\labelenumi{\alph{enumi})}
\setcounter{enumi}{1}
\item
  X 上存在这样的结构，当且仅当 H 在 V 上的伴随表示的诸不可约子表示的交换域都同构于 C；此时这种结构共有 \(2^{s}\) 个，其中 s 是 V 的不可约子表示的个数（利用习题 7）。
\item
  设 j 是 X 上在 G 下不变的殆复结构，\(\mathfrak{p} = \mathfrak{q}(j)\) 是相联的子空间。则 j 可积（即对应于 X 上的一个复解析流形结构，参见~《微分流形与解析流形，结果》，8.8.5 至 8.8.8），当且仅当 p 满足条件
\end{enumerate}

\[
[ \mathfrak {p}, \mathfrak {p} ] \subset \mathfrak {p}.\tag{4}
\]

\begin{enumerate}
\def\labelenumi{\alph{enumi})}
\setcounter{enumi}{3}
\item
  X 上存在在 G 下不变的复结构（即可积的殆复结构），当且仅当 H 是 G 中某个环面的中心化子（证明上面的条件 (1) 至 (4) 蕴含 p 是 \(g_{C}\) 的一个抛物子代数（第 VIII 章，§3，节 5））；此时这些复结构与 \(g_{C}\) 中那些是 \(h_{C}\) 与其幂幺根基的直和的抛物子代数 p 一一对应（参见~第 VIII 章，§3，节 4）。
\item
  证明 G/T 上在 G 下不变的复结构恰有 \(\operatorname{Card}(W)\) 个；若 \(\sigma\) 与 \(\sigma'\) 是两个这样的结构，则存在唯一的元素 \(w \in W\)，使 w 在 G/T 上的典则作用（通过内自同构）把 \(\sigma\) 变为 \(\sigma'\)。若 w 是 W 的非单位元素，则对 G/T 上任何在 G 下不变的复结构，w 在 G/T 上的作用都不是 C-解析的。
\item
  确定 \(S^{2}\) 上在 \(\mathbf{SO}(3)\) 下不变的复结构。
\end{enumerate}

\(g)^{*}\) 采用习题 8 \(d\)) 的记号，设 \(\mathrm{G}_c\) 是 \(\mathrm{G}\) 的复化，\(\mathrm{P}\) 是 \(\mathrm{G}_c\) 中以 \(\mathfrak{p}\) 为李代数的复李子群。证明典则映射 \(\mathrm{G / H} \to \mathrm{G}_c / \mathrm{P}\) 是复解析流形的一个同构.*

\begin{enumerate}
\def\labelenumi{\arabic{enumi})}
\setcounter{enumi}{8}
\tightlist
\item
  设 H 是 G 的一个连通闭子群，它具有极大秩、异于 G，且对这些性质而言是极大的。用 Z 表示商群 \(\mathrm{C}(\mathrm{H})/\mathrm{C}(\mathrm{G})\)。
\end{enumerate}

\begin{enumerate}
\def\labelenumi{\alph{enumi})}
\item
  我们有 \(\dim Z \leq 1\)；若 \(\dim Z = 0\)，则 Z 的阶为 2、3 或 5（化归到 G 概单且中心平凡的情形；应用第 VI 章，§4，习题 4）。
\item
  假定 G 是概单的且 Z 的阶为 3 或 5；置 \(z = \text{Card}(Z)\) 与 \(\pi = \text{Card}(\pi_1(H))/\text{Card}(\pi_1(G))\)。证明我们必定处于下列七种情形之一：
\end{enumerate}

\begin{enumerate}
\def\labelenumi{(\roman{enumi})}
\item
  G 是 \(\mathrm{G}_2\) 型的，H 是 \(\mathrm{A}_2\) 型的，且 \(z = 3, \pi = 1\)；
\item
  G 是 \(F_{4}\) 型的，H 是 \(A_{2} \times A_{2}\) 型的，且 \(z = \pi = 3\)；
\item
  G 是 \(E_{6}\) 型的，H 是 \(A_{2} \times A_{2} \times A_{2}\) 型的，且 \(z = \pi = 3\)；
\item
  G 是 \(E_{7}\) 型的，H 是 \(A_{2} \times A_{5}\) 型的，且 \(z = \pi = 3\)；
\item
  G 是 \(E_{8}\) 型的，H 是 \(A_{8}\) 型的，且 \(z = \pi = 3\)；
\item
  G 是 \(E_{8}\) 型的，H 是 \(A_{2} \times E_{6}\) 型的，且 \(z = \pi = 3\)；
\item
  G 是 \(E_{8}\) 型的，H 是 \(A_{4} \times A_{4}\) 型的，且 \(z = \pi = 5\)。
\end{enumerate}

（利用同上以及第 VI 章的图版；为计算 \(\pi\)，注意若 \(f\) 与 \(f'\) 分别表示 G 与 H 的连通指数，则 \(z\pi f = f'\)。）

\begin{enumerate}
\def\labelenumi{\alph{enumi})}
\setcounter{enumi}{2}
\tightlist
\item
  在上述每种情形中，确定群 H。
\end{enumerate}

\begin{enumerate}
\def\labelenumi{\arabic{enumi})}
\setcounter{enumi}{9}
\tightlist
\item
  我们沿用上一习题的记号。
\end{enumerate}

\begin{enumerate}
\def\labelenumi{\alph{enumi})}
\item
  假定 \(\dim Z = 1\)。则 G/H 上在 G 下不变的复结构恰有两个；存在 G 的一个自同构，它使 H 稳定并互换这两个结构（利用习题 8）。
\item
  确定 \(\mathbf{P}_n(\mathbf{C})\) 上在 \(\mathbf{SU}(n + 1,\mathbf{C})\) 下不变的复结构。
\item
  假定 \(\dim Z = 0\) 且 \(\operatorname{Card}(Z) \neq 2\)。证明 G/H 上存在在 G 下不变的殆复结构（若 z 是 C(H) 中在 G 内不是中心的元素，则 Int z 诱导出 G/H 的一个奇数阶自同构（习题 9）；利用习题 8 b)）。证明 G/H 上不存在在 G 下不变的复结构（利用习题 8 d)）。
\end{enumerate}

\(d\) ) \(\mathbf{S}_6\) 上不存在在 \(\mathbf{SO}(7,\mathbf{R})\) 下不变的复结构（利用 §3 的习题 7）。

\begin{enumerate}
\def\labelenumi{\arabic{enumi})}
\setcounter{enumi}{10}
\tightlist
\item
  我们沿用习题 9 与 10 的记号。
\end{enumerate}

\begin{enumerate}
\def\labelenumi{\alph{enumi})}
\item
  齐性空间 G/H 是对称的（§1，习题 8），当且仅当 Z 的阶为 2 或维数为 1。此时对称空间 G/H 是不可约的（同上）。
\item
  假定 \(\dim Z = 1\)；用 \(X\) 表示复解析流形 \(G / H\)。证明此时我们处于下列情形之一：
\end{enumerate}

\begin{enumerate}
\def\labelenumi{(\roman{enumi})}
\item
  \(\mathbf{G}\) 是 \(\mathrm{A}_l\) 型的，\(\mathrm{D(H)}\) 是 \(\mathrm{A}_{p - 1}\times \mathrm{A}_{l - p}\) 型的；\(\mathbf{X}\) 同构于格拉斯曼流形 \(\mathbf{G}_p(\mathbf{C}^{l + 1})\)。
\item
  G 是 \(B_{l}\) 型的，D(H) 是 \(B_{l-1}\) 型的；X 同构于 \(\mathbf{P}_{2l}(\mathbf{C})\) 中由一个分离二次型的取零所定义的子流形（光滑射影二次曲面）。
\end{enumerate}

(ii') G 是 \(D_{l}\) 型的，\(D(H)\) 是 \(D_{l-1}\) 型的；X 同构于 \(\mathbf{P}_{2l-1}(\mathbf{C})\) 中的一个光滑射影二次曲面。

\begin{enumerate}
\def\labelenumi{(\roman{enumi})}
\setcounter{enumi}{2}
\item
  G 是 \(C_{l}\) 型的，D(H) 是 \(A_{l-1}\) 型的；X 同构于 \(\mathbf{G}_{l}(\mathbf{C}^{2l})\) 中由 \(C^{2l}\) 上通常的交错双线性型的极大迷向子空间组成的子流形。
\item
  G 是 \(D_{l}\) 型的，\(\mathrm{D}(\mathrm{H})\) 是 \(A_{l-1}\) 型的；X 同构于 \(\mathbf{G}_{l}(\mathbf{C}^{2l})\) 中由 \(C^{2l}\) 上通常的对称双线性型的极大迷向子空间组成的子流形。
\item
  G 是 \(E_{6}\) 型的，D(H) 是 \(D_{5}\) 型的。
\item
  G 是 \(E_{7}\) 型的，D(H) 是 \(E_{6}\) 型的。
\end{enumerate}

\begin{enumerate}
\def\labelenumi{\alph{enumi})}
\setcounter{enumi}{2}
\tightlist
\item
  假定 \(\operatorname{Card}(Z)=2\)；给出可能情形的列表。若 G 是 \(A_{l}\)、\(B_{l}\)、\(C_{l}\) 或 \(D_{l}\) 型的，证明实流形 G/H 同构于一个格拉斯曼流形 \(\mathrm{G}_{p}(\mathrm{K}^{q})\)，其中 K=R、C 或 H。
\end{enumerate}

¶ 12) 假定 G 是单连通的；对所有 \(\alpha \in \mathrm{R}(\mathrm{G},\mathrm{T})\)，置 \(t_{\alpha} = \exp (\frac{1}{2} K_{\alpha})\)。置 \(\mathrm{N} = \mathrm{N}_{\mathrm{G}}(\mathrm{T})\)，并用 \(\varphi : \mathrm{N} \to \mathrm{W}\) 表示典则映射。设 B 是 \(\mathrm{R}(\mathrm{G},\mathrm{T})\) 的一个基；选取元素 \(n_{\alpha}\)（在 \((\mathrm{N} \cap \mathrm{S}_{\alpha}) - (\mathrm{T} \cap \mathrm{S}_{\alpha})\) 中，对所有 \(\alpha \in \mathrm{B}\)）。

\begin{enumerate}
\def\labelenumi{\alph{enumi})}
\item
  证明 \(\varphi(n_{\alpha}) = s_{\alpha}\) 且 \(n_{\alpha}^{2} = t_{\alpha}\)，从而 \(n_{\alpha}^{4} = 1\)。
\item
  设 \(\alpha\) 与 \(\beta\) 是 B 的两个不同元素，并设 \(m_{\alpha \beta}\) 是 \(s_{\alpha}s_{\beta}\) 在 W 中的阶。证明
\end{enumerate}

\[
\begin{array}{r l} {n _ {\alpha} n _ {\beta} = n _ {\beta} n _ {\alpha}} & {\mathrm{if} m _ {\alpha \beta} = 2} \\ {n _ {\alpha} n _ {\beta} n _ {\alpha} = n _ {\beta} n _ {\alpha} n _ {\beta}} & {\mathrm{if} m _ {\alpha \beta} = 3} \\ {(n _ {\alpha} n _ {\beta}) ^ {2} = (n _ {\beta} n _ {\alpha}) ^ {2}} & {\mathrm{if} m _ {\alpha \beta} = 4} \\ {(n _ {\alpha} n _ {\beta}) ^ {3} = (n _ {\beta} n _ {\alpha}) ^ {3}} & {\mathrm{if} m _ {\alpha \beta} = 6} \end{array}
\]

（例如若 \(m_{\alpha\beta}=3\)，则 \((s_{\alpha}s_{\beta})s_{\alpha}(s_{\alpha}s_{\beta})^{-1}=s_{\beta}\) 且 \(s_{\alpha}s_{\beta}(\alpha)=\beta\)；证明 \(n_{\alpha}n_{\beta}n_{\alpha}n_{\beta}^{-1}n_{\alpha}^{-1}n_{\beta}^{-1}\) 属于 \(S_{\beta}\)，并由 \(S_{\alpha}\cap S_{\beta}\cap T=\{e\}\) 得出结论）。

\begin{enumerate}
\def\labelenumi{\alph{enumi})}
\setcounter{enumi}{2}
\item
  由 b) 推出，存在唯一的截面 \(\nu: W \to N\)（\(\varphi\) 的），使得 \(\nu(s_{\alpha}) = n_{\alpha}\)，且 \(\nu(ww') = \nu(w)\nu(w')\)（当 \(l(ww') = l(w) + l(w')\) 时）（其中 \(l(w)\) 表示 w 关于生成系 \((s_{\alpha})_{\alpha \in \mathbb{B}}\) 的长度；参见~第 VI 章，§1，节 5，命题 5）。置 \(n_w = \nu(w)\)。
\item
  设 \(W^{*}\) 是 N 中由诸 \(n_{\alpha}\) 生成的子群；证明 \(W^{*} \cap T\) 是子群 \(T_{2}\)（T 中由阶 \(\leq 2\) 的元素组成），且 W 可与 \(W^{*}/T_{2}\) 等同。
\item
  设 \(w \in \mathrm{W}\) 满足 \(w^2 = 1\)；证明 \(n_w^2 = \prod_{\alpha \in \mathrm{R}_w} t_\alpha\)，其中 \(\mathrm{R}_w\) 是正根 \(\alpha\)（满足 \(w(\alpha) < 0\)）组成的集合（写 \(w = s_{\alpha_1} \ldots s_{\alpha_r}\)，其中 \(r = l(w)\) 且 \(\alpha_i \in \mathrm{B}\)；应用 c) 以及第 VI 章，§1，节 6，命题 17 的推论 2）。
\end{enumerate}

\(f\) ) 假定 \(\mathrm{G}\) 是概单的。设 \(c\) 是 \(\mathrm{W}\) 中的一个 Coxeter 变换，\(h\) 是 \(\mathrm{W}\) 的 Coxeter 数（第 VI 章，§1，节 11）。证明 \(n_c^h = \prod_{\alpha \in \mathbb{R}_+}t_\alpha\)（利用 \(c\)）、\(e\)）以及第 V 章，§6 的习题 2）。

\begin{enumerate}
\def\labelenumi{\arabic{enumi})}
\setcounter{enumi}{12}
\tightlist
\item
  \begin{enumerate}
  \def\labelenumii{\alph{enumii})}
  \tightlist
  \item
    选取 R(G, T) 的一个基 B，并用 R+ 表示 R(G, T) 的正根集合。证明 \(z_{G} = \prod_{\alpha \in R_{+}} \exp(\frac{1}{2} K_{\alpha})\) 是 C(G) 的一个元素，它与 T 和 B 的选取无关；我们有 \(z_{G}^{2} = e\)。
  \end{enumerate}
\end{enumerate}

\begin{enumerate}
\def\labelenumi{\alph{enumi})}
\setcounter{enumi}{1}
\item
  设 H 是另一个连通紧李群。则 \(z_{\mathrm{G}\times\mathrm{H}} = (z_{\mathrm{G}}, z_{\mathrm{H}})\)；证明若 \(f : G \to H\) 是李群的满射态射，则 \(f(z_{\mathrm{G}}) = z_{\mathrm{H}}\)。
\item
  置 \(\mathrm{R} = \mathrm{R}(\mathrm{G},\mathrm{T})\)；假定 \(\mathbf{G}\) 是单连通的，从而 \(\mathrm{X(T)}\) 可与 \(\mathrm{P(R)}\) 等同。证明同态 \(\chi \mapsto \chi (z_{\mathrm{G}})\)
\end{enumerate}

（从 X(T) 到 \(\{1,-1\}\)）的核是第 VI 章，§1，习题 8 中定义的子群 \(P'(R)\)。

\begin{enumerate}
\def\labelenumi{\alph{enumi})}
\setcounter{enumi}{3}
\tightlist
\item
  若 \(\mathrm{G} = \mathbf{S}\mathbf{U}(n,\mathbf{C})\)，则 \(z_{\mathrm{G}} = (-1)^{n + 1}I_n\)；
\end{enumerate}

若 \(\mathrm{G} = \mathbf{S}\mathbf{U}(n,\mathbf{H})\)，则 \(z_{\mathrm{G}} = -I_n\)；

若 \(\mathrm{G} = \mathbf{Spin}(n, \mathbf{R})\) 且 \(n \equiv 3, 4, 5\) 或 6 (mod 8)，则 \(z_{G} = -1\)；

若 \(\mathrm{G} = \mathbf{Spin}(n, \mathbf{R})\) 且 \(n \equiv 0, 1, 2, 7 \pmod{8}\)，则 \(z_{G} = 1\)；

若 G 是 \(E_{6}\)、\(E_{8}\)、\(F_{4}\) 或 \(G_{2}\) 型的，则 \(z_{G} = e\)；

若 G 是 \(E_{7}\) 型单连通的，则 \(z_{G}\) 是 \(\mathrm{C}(\mathrm{G})\) 的唯一的非单位元素。

（利用第 VI 章，§4，习题 5。）

\begin{enumerate}
\def\labelenumi{\arabic{enumi})}
\setcounter{enumi}{13}
\tightlist
\item
  假定群 G 是概单的；用 h 表示 R(G,T) 的 Coxeter 数（第 VI 章，§1，节 11）。若存在 G 的一个极大环面 S，使 g 属于 \(\mathrm{N}_{\mathrm{G}}(\mathrm{S})\)，且 g 在 \(\mathrm{W}_{\mathrm{G}}(\mathrm{S})\) 中的类是一个 Coxeter 变换（同上），则称 G 的元素 g 为 Coxeter 元。
\end{enumerate}

\begin{enumerate}
\def\labelenumi{\alph{enumi})}
\item
  证明任意两个 Coxeter 元都共轭（按节 5 命题 10 的推论的证明方式论证）。
\item
  每个 Coxeter 元 g 都是正则的，且满足 \(g^{h} = z_{G}\)，其中 \(z_{G}\) 是习题 13 中定义的 C(G) 的元素；特别地，视 \(z_{G}\) 等于或不等于 e，g 的阶为 h 或 2h（利用习题 12 f)）。
\item
  元素 \(g \in G\) 是 Coxeter 元，当且仅当 \(Ad g \otimes 1_{C}\)（\(g_{C}\) 的自同构）满足第 VIII 章，§5，习题 5 f) 的诸等价条件。
\end{enumerate}

\(d\) ) 证明每个正则元 \(g\)（属于 \(\mathbf{G}\)，满足 \(g^{h}\in \mathrm{C}(\mathrm{G})\)）都是 Coxeter 元；当 \(p < h\) 时，不存在正则元 \(k\) 满足 \(k^p\in \mathrm{C}(\mathrm{G})\)。

\begin{enumerate}
\def\labelenumi{\arabic{enumi})}
\setcounter{enumi}{14}
\tightlist
\item
  设 H 是 G 的一个连通闭子群。若 H 不包含于任何异于 G 的极大秩连通闭子群之中，则称 H 是净的。
\end{enumerate}

\begin{enumerate}
\def\labelenumi{\alph{enumi})}
\item
  证明 H 是净的，当且仅当它在 G 中的中心化子等于 C(G)。特别地，若 H 是净的，则 C(H) = C(G) ∩ H。
\item
  以下假定 H 是净的。证明对 H 的每个极大环面 S，都有 \(\mathrm{C}(\mathrm{H}) = \mathrm{S} \cap \mathrm{C}(\mathrm{G})\)。
\item
  设 \(H'\) 是 G 的包含 H 的一个连通闭子群。证明 \(rkH < rkH'\)，并由此推出 H 在 \(H'\) 中是净的。
\item
  设 K 是 H 的一个连通闭子群，它在 H 中是净的且包含 G 的一个正则元。证明 K 在 G 中是净的。
\end{enumerate}

\begin{enumerate}
\def\labelenumi{\arabic{enumi})}
\setcounter{enumi}{15}
\tightlist
\item
  设 H 是 G 的一个连通闭子群，使得 \(T \cap H\) 是 H 的一个极大环面 S。设 \(\lambda \in \mathrm{R}(\mathrm{H}, \mathrm{S})\)；用 \(\mathrm{R}(\lambda)\) 表示 \(\mathrm{R}(\mathrm{G}, \mathrm{T})\) 中在 S 上的限制等于 \(\lambda\) 的根组成的集合。
\end{enumerate}

\begin{enumerate}
\def\labelenumi{\alph{enumi})}
\item
  证明 \(\mathrm{R}(\lambda)\) 非空。
\item
  设 \(w \in \mathrm{W}_{\mathrm{H}}(\mathrm{S})\)；证明存在元素 \(\bar{w}\)（属于 \(\mathrm{W}_{\mathrm{G}}(\mathrm{T})\)），使 \(\mathrm{R}(w\lambda) = \bar{w}\mathrm{R}(\lambda)\)（利用 §2 的习题 4）。
\item
  设 \(\mathrm{P}(\lambda)\) 是 \(\mathrm{R}(\mathrm{G},\mathrm{T})\) 与 \(\mathrm{X}(\mathrm{T})\) 中由 \(\mathrm{R}(\lambda)\) 生成的子群之交。证明存在 G 的包含 T 的闭子群 \(G_{\lambda}\)，使 \(\mathrm{P}(\lambda)=\mathrm{R}(\mathrm{G}_{\lambda},\mathrm{T})\)。由此推出反射 \(s_{\lambda}\in\mathrm{W}_{\mathrm{H}}(\mathrm{S})\) 是诸反射 \(s_{\alpha}\)（\(\alpha\in\mathrm{P}(\lambda)\)）之积在 S 上的限制。
\item
  证明结点向量 \(K_{\lambda} \in \mathrm{L}(\mathrm{S})\)（对应于 \(\lambda\)）是诸 \(K_{\alpha}\)（\(\alpha \in \mathrm{R}(\lambda)\)）的一个整系数线性组合。
\item
  设 \(B_{H}\) 是 \(R(H,S)\) 的一个基。证明 \(R(\lambda)\) 包含于 \(X(T)\) 中由诸 \(R(\mu)\)（\(\mu \in B_{H}\)）的并所生成的子群（证明具有所述性质的 \(\lambda \in R(H,S)\) 组成的集合在所有 \(s_{\mu}\)（\(\mu \in B_{H}\)）下稳定，并利用 c)）。
\end{enumerate}

¶ 17) 我们沿用上一习题的记号；再假定子群 H 是净的（习题 15）。

\begin{enumerate}
\def\labelenumi{\alph{enumi})}
\item
  证明 \(\mathrm{R}(\mathrm{G},\mathrm{T})\) 包含于 \(\mathrm{X}(\mathrm{T})\) 中由诸 \(\mathrm{R}(\lambda)\)（\(\lambda\in B_{H}\)）的并所生成的子群。
\item
  设 \(\Delta\) 是 S 中由诸 \(s \in S\)（满足 \(\lambda(s) = \mu(s)\)，对所有 \(\lambda, \mu\) 属于 \(B_{H}\)）组成的子群的单位元连通分支。证明存在 R(G,T) 的一个基 \(\{\alpha_{1}, \ldots, \alpha_{l}\}\) 和一个整数 k，\(0 \leq k \leq l - 1\)，使 \(\Delta\) 是满足下列条件的 \(t \in T\) 组成的集合的单位元连通分支
\end{enumerate}

\[
\alpha_ {1} (t) = \dots = \alpha_ {k} (t) = 1, \alpha_ {k + 1} (t) = \dots = \alpha_ {l} (t).
\]

（设 \(x\) 是 \(\mathrm{L}(\mathrm{S})\) 中使 \(\delta(\lambda)(x) = 2\pi i\)（对所有 \(\lambda \in \mathrm{B}_{\mathrm{H}}\)）成立的元素；由 \(a\)) 推出 \(\exp x \in \mathrm{C}(\mathrm{G})\)。选取 \(\{\alpha_1, \ldots, \alpha_l\}\)，使 \(i\delta(\alpha_j)(x)\) 为零（对 \(1 \leq j \leq k\)）且 \(< 0\)（对 \(k + 1 \leq j \leq l\)）；然后证明对所有 \(\lambda \in \mathrm{B}_{\mathrm{H}}\)，\(\mathrm{R}(\lambda)\) 中的每个根都可写成

\[
\alpha_ {j} + n _ {1} \alpha_ {1} + \dots + n _ {k} \alpha_ {k},
\]

其中 \(j \geq k + 1\) 且 \(n_i \in \mathbf{N}\)。利用 \(a)\) 得出结论。

\begin{enumerate}
\def\labelenumi{\alph{enumi})}
\setcounter{enumi}{2}
\tightlist
\item
  证明整数 k 不依赖于环面 S、T 或基 \(B_{H}, \{\alpha_{1}, \ldots, \alpha_{l}\}\) 的选取；它等于 \(Z_{\mathrm{G}}(\Delta)\) 的换位子群的秩。
\end{enumerate}

\begin{enumerate}
\def\labelenumi{\arabic{enumi})}
\setcounter{enumi}{17}
\tightlist
\item
  我们沿用习题 16 与 17 的记号。若子群 H 是净的，且子环面 \(\Delta\) 包含一个正则元（换言之，k = 0），则称子群 H 是主的。
\end{enumerate}

\begin{enumerate}
\def\labelenumi{\alph{enumi})}
\item
  设 \(\mathrm{K}\) 是 \(\mathrm{H}\) 的一个连通闭子群。证明 \(\mathrm{K}\) 是 \(\mathrm{G}\) 的主子群，当且仅当 \(\mathrm{K}\) 在 \(\mathrm{H}\) 中是主的且 \(\mathrm{H}\) 在 \(\mathrm{G}\) 中是主的（利用习题 15 c) 与 d)）。
\item
  以下假定 H 是主的。证明诸 \(\mathrm{R}(\mu)\)（\(\mu \in B_{H}\)）的并是 \(\mathrm{R}(\mathrm{G}, \mathrm{T})\) 的一个基。
\item
  设 \(\alpha \in \mathrm{R}(\mathrm{G},\mathrm{T})\)。证明 \(\alpha\) 在 S 上的限制是某个根 \(\lambda \in \mathrm{R}(\mathrm{H},\mathrm{S})\) 的非零整数倍；根 \(\alpha\) 是 \(\mathrm{R}(\lambda)\) 中元素之和（证明 \(\mathrm{L}(\mathrm{S})\) 与超平面 \(\delta (\alpha) = 0\) 的交是 \(\mathrm{L}(\mathrm{S}))\) 的一面墙）。
\item
  设 \(\lambda\in\mathrm{R}(\mathrm{H},\mathrm{S})\)；证明结点向量 \(K_{\lambda}\) 是诸 \(K_{\alpha}\)（\(\alpha\in\mathrm{R}(\lambda)\)）的一个非负整系数线性组合（参见~习题 16 d) 以及第 V 章，§3，节 5，引理 6）。
\end{enumerate}

\begin{enumerate}
\def\labelenumi{\arabic{enumi})}
\setcounter{enumi}{18}
\tightlist
\item
  我们沿用习题 16 至 18 的记号；假定子群 H 是主的。给 \(\mathrm{X(T)\otimes\mathbf{R}}\)（相应地 \(\mathrm{X(S)\otimes\mathbf{R}}\)）赋以一个在 \(\mathrm{W_{G}(T)}\)（相应地 \(\mathrm{W_{H}(S)}\)）下不变的内积。设 \(\lambda,\mu\) 是 \(B_{H}\) 中的两个根。
\end{enumerate}

\begin{enumerate}
\def\labelenumi{\alph{enumi})}
\item
  证明若 \(\lambda\) 与 \(\mu\) 正交，则集合 \(\mathrm{R}(\lambda)\) 与 \(\mathrm{R}(\mu)\) 正交。由此推出，若 \(\mathrm{G}\) 是概单的，则 \(\mathrm{H}\) 也是概单的。
\item
  以下假定 \(n(\lambda, \mu) = -1\)。证明存在一个满射 \(u: \mathrm{R}(\lambda) \to \mathrm{R}(\mu)\)，使得对 \(\alpha \in \mathrm{R}(\lambda)\)，\(u(\alpha)\) 是 \(\mathrm{R}(\mu)\) 中与 \(\alpha\) 相连（即不正交）的唯一根；我们有 \(K_{\mu} = \sum_{\beta \in \mathrm{R}(\mu)} K_{\beta}\)，且 \(\mathrm{R}(\mu)\) 中的根两两正交（把 \(K_{\mu}\) 与 \(K_{\lambda}\) 写成诸 \(K_{\alpha}\) 的线性组合，然后比较系数）。
\item
  若 \(n(\mu, \lambda) = -1\)，则映射 \(u\) 是双射，且 \(\mathrm{R}(\lambda) \cup \mathrm{R}(\mu)\) 中相连的根对只有诸对 \((\alpha, u(\alpha))\)（\(\alpha \in \mathrm{R}(\lambda)\)）。
\item
  假定 \(n(\mu,\lambda)=-2\)；设 \(\beta\in\mathrm{R}(\mu)\)。证明 \(u^{-1}(\beta)\) 含有一个或两个元素；若 \(u^{-1}(\beta)=\{\alpha_{1},\alpha_{2}\}\)，则根 \(\alpha_{i}\)（i=1,2）与 \(\mathrm{R}(\lambda)\) 中其余的根正交，且与 \(\beta\) 等长；若 \(u^{-1}(\beta)=\{\alpha\}\) 且 \(\|\alpha\|=\|\beta\|\)，则 \(\alpha\) 与唯一的根 \(\alpha^{\prime}\in\mathrm{R}(\lambda)\) 相连，后者与 \(\alpha\) 等长，且 \(\{\alpha,\alpha^{\prime}\}\) 与 \(\mathrm{R}(\lambda)\) 的其余部分正交；若 \(u^{-1}(\beta)=\{\alpha\}\) 且 \(\|\alpha\|\neq\|\beta\|\)，则 \(\alpha\) 与 \(\mathrm{R}(\lambda)\) 中其余的根正交。
\item
  类似地讨论 \(n(\mu, \lambda) = -3\) 的情形。
\end{enumerate}

\begin{enumerate}
\def\labelenumi{\arabic{enumi})}
\setcounter{enumi}{19}
\tightlist
\item
  假定群 G 是概单的；设 H 是 G 的一个秩 \(\geq 2\) 的主子群。
\end{enumerate}

\begin{enumerate}
\def\labelenumi{\alph{enumi})}
\item
  证明 H 是 \(B_{h}, C_{h}, F_{4}\) 或 \(G_{2}\) 型的半单群（利用习题 19）。
\item
  假定 H 的秩 \(\geq 3\)。证明 G 是 \(A_{l}, D_{l}, E_{6}, E_{7}\) 或 \(E_{8}\) 型的（考察 R(G, T) 的邓金图的端点，并应用习题 19）。
\item
  证明若 \(rkH \geq 3\)，则我们处于下列情形之一：
\end{enumerate}

G 是 \(A_{2l}\) 型 \((l \geq 3)\) 的，H 是 \(B_{l}\) 型的；

G 是 \(\mathrm{A}_{2l - 1}\) 型 \((l\geq 3)\) 的，H 是 \(\mathbf{C}_l\) 型的

G 是 \(D_{l}\) 型 \((l \geq 4)\) 的，H 是 \(B_{l-1}\) 型的；

G 是 \(\mathrm{E}_6\) 型的，H 是 \(\mathrm{F}_4\) 型的。

\(d\) ) 证明若 H 是 \(\mathrm{B}_2\) 型的，则 G 是 \(\mathrm{A}_3\) 或 \(\mathrm{A}_4\) 型的。

\begin{enumerate}
\def\labelenumi{\alph{enumi})}
\setcounter{enumi}{4}
\tightlist
\item
  证明若 H 是 \(G_{2}\) 型的，则 G 是 \(B_{3}, D_{4}\) 或 \(A_{6}\) 型的。
\end{enumerate}

（关于这些情形更明确的描述，参见 §5，习题 5。）

\(f)\) 设 \(\mathrm{K}\) 是 \(\mathrm{G}\) 的包含 \(\mathrm{H}\) 的一个连通闭子群；证明或者 \(\mathrm{K} = \mathrm{G}\)，或者 \(\mathrm{K} = \mathrm{H}\)，或者 \(\mathrm{H}\) 是 \(\mathrm{G}_2\) 型的、\(\mathrm{K}\) 是 \(\mathrm{B}_3\) 型的且 \(\mathrm{G}\) 是 \(\mathrm{D}_4\) 或 \(\mathrm{A}_6\) 型的。

\begin{enumerate}
\def\labelenumi{\arabic{enumi})}
\setcounter{enumi}{20}
\tightlist
\item
  \begin{enumerate}
  \def\labelenumii{\alph{enumii})}
  \tightlist
  \item
    设 H 是 G 的一个秩为 1 的连通闭子群。证明下列条件等价：
  \end{enumerate}
\end{enumerate}

\begin{enumerate}
\def\labelenumi{(\roman{enumi})}
\item
  H 是主的；
\item
  H 是净的且包含 G 的一个正则元；
\item
  存在一个主 \(\mathfrak{sl}_2\)-三元组 \((x,h,y)\)，它位于 \(\mathfrak{g}_{\mathbf{C}}\) 中（第 VIII 章，§11，节 4），使得
\end{enumerate}

\[
\mathrm{L} (\mathrm{H}) _ {(\mathbf {C})} = \mathbf {C} x + \mathbf {C} h + \mathbf {C} y.
\]

\begin{enumerate}
\def\labelenumi{\alph{enumi})}
\setcounter{enumi}{1}
\item
  证明 G 包含一个秩为 1 的主子群，且任意两个这样的子群都共轭（用习题 17 的记号，注意 \(S = \Delta\) ）。
\item
  证明 G 的一个连通闭子群是主的，当且仅当它包含 G 的一个秩为 1 的主子群（利用习题 18 a)）。
\end{enumerate}

\begin{enumerate}
\def\labelenumi{\arabic{enumi})}
\setcounter{enumi}{21}
\item
  设 H 是 G 的一个秩为 1 的主子群；设 \(\Gamma\) 是 \(\operatorname{Aut}(G)\) 中由满足 \(u(\mathrm{H}) = \mathrm{H}\) 的自同构 u 组成的子群。证明 \(\operatorname{Aut}(G) = \Gamma.\operatorname{Int}(G)\)。
\end{enumerate}

\exercisesection{\S~5}

\begin{enumerate}
\def\labelenumi{\arabic{enumi})}
\item
  假定 G 是单连通的；G 的单房是 G 中形如 \(\exp(A)\) 的子集，其中 A 是 g 的某个嘉当子代数的一个单房。

\begin{enumerate}
\def\labelenumi{\alph{enumi})}
\item
  证明 \(\mathbf{G}\) 的诸单房构成 \(\mathrm{G}_r\) 的一个分划。
\item
  G 的每个单房都包含于唯一的极大环面。
\item
  证明 G 的包含于 T 的诸单房构成 \(T_{r}\) 的一个分划，且这些单房组成的集合是 W 下的一个主齐性空间。
\item
  G 的正则元素的每个共轭类在每个单房中都有唯一的点。
\item
  设 E 是 G 的一个单房。证明 E 是一个可缩空间；若 g 是单的，则 E 同胚于一个欧几里得空间中的开单纯形。
\end{enumerate}

\begin{enumerate}
\def\labelenumi{\arabic{enumi})}
\setcounter{enumi}{1}
\tightlist
\item
  \begin{enumerate}
  \def\labelenumii{\alph{enumii})}
  \tightlist
  \item
    设 E 是一个单连通拓扑空间，\(u : E \to G\) 是满足 \(u(E) \subset G_r\) 的连续映射。证明 u（《一般拓扑学》第 XI 章）同伦于以 e 为值的常值映射（考察覆盖 \(\varphi_r : (G/T) \times t_r \to G_r\)；把 u 提升为 \(\tilde{u} : E \to (G/T) \times t\)，再利用 t 可缩这一事实）。
  \end{enumerate}
\end{enumerate}

\(^{*}b)\) 证明群 \(\pi_{2}(G)\) 为零。

（设 \(u: \mathbf{S}_2 \to \mathbf{G}\) 是一个 \(\mathrm{C}^\infty\) 类映射；利用命题 1 与横截性定理，证明 \(u\) 同伦于一个像包含于 \(\mathrm{G}_r\) 中的映射，然后应用 \(a).)_*\)

\begin{enumerate}
\def\labelenumi{\arabic{enumi})}
\setcounter{enumi}{2}
\tightlist
\item
  设 \(f: \tilde{\mathrm{G}} \to \mathrm{G}\) 是 \(\mathrm{G}\) 的一个泛覆盖，\(\pi\) 是 \(f\) 的核（同构于 \(\pi_1(\mathrm{G})\)）；置 \(\mathrm{C} = \mathrm{C}(\mathrm{G})\) 与 \(\tilde{\mathrm{C}} = \mathrm{C}(\tilde{\mathrm{G}})\)。设 \(\sigma\) 是
\end{enumerate}

G 的一个自同构，\(\tilde{\sigma}\) 是 \(\tilde{G}\) 的一个自同构（由 \(\sigma\) 诱导）。用 \(G_{\sigma}, C_{\sigma}, \tilde{G}_{\sigma}, \tilde{C}_{\sigma}\) 分别表示 G, C, \(\tilde{G}\)，\(\tilde{C}\) 中被 \(\sigma, \sigma, \tilde{\sigma}, \tilde{\sigma}\) 固定的点组成的集合。

\begin{enumerate}
\def\labelenumi{\alph{enumi})}
\item
  证明 \((\mathrm{G}_{\sigma})_0\)（\(\mathrm{G}_{\sigma}\) 的单位元连通分支）是 \(f(\tilde{\mathrm{G}}_{\sigma})\)。
\item
  用 s 表示 Z-模 \(\pi\) 的由 \(\tilde{\sigma}\) 诱导的自同态。证明商群 \(\mathrm{G}_{\sigma}/(\mathrm{G}_{\sigma})_{0}\) 同构于 \(\operatorname{Coker}(1-s)\) 的一个子群，特别地它是交换的。若 1-s 是满射，则 \(G_{\sigma}\) 是连通的。
\item
  设 \(n \in N\) 满足 \(s^{n} = \operatorname{Id}_{\pi}\)。证明群 \(\mathrm{G}_{\sigma}/(\mathrm{G}_{\sigma})_{0}\) 可与商 \(\operatorname{Ker}(1 + s + \cdots + s^{n-1})/\operatorname{Im}(1 - s)\) 的一个子群等同；由此推出它被 n 零化。~若 n 与 \(\pi\) 的挠子群的阶互素，则 \(G_{\sigma}\) 是连通的。
\end{enumerate}

\(d\) ) 利用《代数学》（Algèbre）第 X 章第~194 页的习题 23，重新得出 \(b\)) 与 \(c\)) 的结果。

\begin{enumerate}
\def\labelenumi{\alph{enumi})}
\setcounter{enumi}{4}
\tightlist
\item
  证明在下列每种情形中 \(\mathrm{G}_{\sigma}\) 都是连通的：
\end{enumerate}

\begin{enumerate}
\def\labelenumi{(\roman{enumi})}
\item
  G 是 \(\mathrm{A}_{2n}\) 型（\(n \geq 1\)）的半单群且 \(\sigma\) 不是内的；
\item
  G 是 \(\mathrm{E}_6\) 型的半单群且 \(\sigma\) 不是内的；
\item
  G 是 \(\mathrm{D}_4\) 型的半单群且 \(\sigma\) 是一个三重性（即模 \(\operatorname{Int}(\mathbf{G})\) 后的阶为 3）。
\end{enumerate}

\(f\) ) 定义一个从 \(\mathrm{C}_{\sigma} / (\mathrm{C}_{\sigma}\cap (\mathrm{G}_{\sigma})_{0})\) 到 \(((1 - \tilde{\sigma})\tilde{\mathrm{C}}\cap \pi) / (1 - \tilde{\sigma})\pi .\) 的同构。由此推出，若 \(1 - s\) 不是满射且 \(\pi \subset (1 - \tilde{\sigma})\mathrm{C}\)，则 \(\mathrm{C}_{\sigma}\) 不包含于 \((\mathrm{G}_{\sigma})_{0}\)。对 \(\mathrm{G} = \mathbf{SO}(2n,\mathbf{R})\) 且 \(\sigma\) 非内时，有 \(-I_{2n}\notin (\mathrm{G}_{\sigma})_{0}\)，且 \(\mathrm{G}_{\sigma}\) 不是连通的。

\begin{enumerate}
\def\labelenumi{\arabic{enumi})}
\setcounter{enumi}{3}
\tightlist
\item
  假定 G 是半单的。设 e 是 G 的一个标架，\(\Phi\) 是 G 的一个尊重该标架的自同构群；用 H 表示 G 中由被 \(\Phi\) 固定的元素组成的子群。
\end{enumerate}

\begin{enumerate}
\def\labelenumi{\alph{enumi})}
\item
  证明 \(H_{0}\) 是半单的；若 G 是单连通的，则 H 是连通的（化归到 G 概单且 \(\Phi\) 循环的情形，并应用定理 1（节 3）以及第 VIII 章，§5，习题 13）。
\item
  假定 \(\mathbf{G}\) 是概单的。证明
\end{enumerate}

若 G 是 \(\mathrm{A}_{2l}\) 型（\(l \geq 1\)）的且 \(\varPhi\) 的阶为 2，则 \(\mathrm{H}_0\) 是 \(\mathrm{B}_l\) 型的；

若 G 是 \(A_{2l-1}\) 型（\(l \geq 2\)）的且 \(\Phi\) 的阶为 2，则 \(H_{0}\) 是 \(C_{l}\) 型的；

若 G 是 \(D_{l}\) 型 \((l \geq 4)\) 的且 \(\Phi\) 的阶为 2，则 \(H_{0}\) 是 \(B_{l-1}\) 型的；

若 G 是 \(D_{4}\) 型的且 \(\Phi\) 的阶为 3 或 6，则 \(H_{0}\) 是 \(G_{2}\) 型的；

若 \(\mathrm{G}\) 是 \(\mathrm{E}_6\) 型的且 \(\varPhi\) 的阶为 2，则 \(\mathrm{H}_0\) 是 \(\mathrm{F}_4\) 型的。

在每种情形中，确定群 \(\pi_i(\mathrm{H})\)，\(i = 0,1\)（参见~习题 3）。

\begin{enumerate}
\def\labelenumi{\alph{enumi})}
\setcounter{enumi}{2}
\tightlist
\item
  证明子群 \(\mathrm{H}_0\) 是主的（§4，习题 18）。
\end{enumerate}

\(d\) ) 证明 \(\mathrm{B}_3\) 或 \(\mathrm{A}_6\) 型的半单群包含一个 \(\mathrm{G}_2\) 型的主子群（利用 \(b\)) 以及 §4 的习题 20）。

\begin{enumerate}
\def\labelenumi{\arabic{enumi})}
\setcounter{enumi}{4}
\tightlist
\item
  假定群 G 是概单的；设 H 是 G 的一个主连通闭子群（§4，习题 18）。用 \(\Phi\) 表示 G 中使 H 固定的自同构 u 组成的群，用 F 表示 G 中被 \(\Phi\) 固定的元素组成的子群。
\end{enumerate}

\begin{enumerate}
\def\labelenumi{\alph{enumi})}
\item
  证明存在 \(\mathrm{G}\) 的一个在 \(\varPhi\) 下稳定的标架。
\item
  证明我们必定处于下列情形之一：
\end{enumerate}

\begin{enumerate}
\def\labelenumi{(\roman{enumi})}
\item
  \(H = F_{0};\)
\item
  G 是 \(B_{3}\) 型的，H 是 \(G_{2}\) 型的，且 \(\Phi\) 化为单位元；
\item
  G 是 \(A_{6}\) 型的，H 是 \(G_{2}\) 型的，且 \(\Phi\) 的阶为 2。
\end{enumerate}

（利用习题 4 以及 §4 的习题 20。）

¶ 6) 假定 G 是单连通的。设 p 是一个素数，g 是 C(G) 中满足 \(g^{p} = e\) 的一个元素。

\begin{enumerate}
\def\labelenumi{\alph{enumi})}
\tightlist
\item
  证明存在元素 \(u \in T\) 与 \(w \in W\)，使得
\end{enumerate}

\begin{enumerate}
\def\labelenumi{(\roman{enumi})}
\item
  \(w(u)u^{-1} = g;\)
\item
  \(w^{p} = 1;\)
\item
  \(u^{p} = e\) 当 \(p \neq 2\)，\(u^{p} = e\) 或 g 当 p = 2。
\end{enumerate}

（设 A 是 t 中的一个单房；对 \(i = 0, 1, \ldots, p - 1\)，设 \(x_i \in \overline{A}\) 满足 \(\exp x_i = g^i\)。取 u 为元素 \(\exp x\)，其中 x 是 \(\overline{A}\) 的以诸 \(x_i\) 为顶点的那个面的重心，并取 w 为 W 中满足 \(w(\mathrm{A}) = \mathrm{A} - x\) 的元素。）

\begin{enumerate}
\def\labelenumi{\alph{enumi})}
\setcounter{enumi}{1}
\tightlist
\item
  证明存在 \(u \in \mathrm{T}\) 与 \(v \in \mathrm{N}_{\mathrm{G}}(\mathrm{T})\)，使得
\end{enumerate}

\begin{enumerate}
\def\labelenumi{(\roman{enumi})}
\item
  \(vuv^{-1}u^{-1}=g;\)
\item
  \(v^{p} = e\) 当 \(p \neq 2\)，\(v^{p} = e\) 或 g 当 p = 2；
\item
  \(u^{p} = e\) 当 \(p \neq 2\)，\(u^{p} = e\) 或 \(g\) 当 \(p = 2\)。
\end{enumerate}

（利用 \(a\)) 中的构造，按 §4 习题 12 的方式把 \(w\) 提升为 \(n_w\)；取 \(v = n_w^{p+1}\) 当 \(p \neq 2\)，取 \(v = n_w\) 当 \(p = 2\)。）

\begin{enumerate}
\def\labelenumi{\arabic{enumi})}
\setcounter{enumi}{6}
\tightlist
\item
  设 p 是一个素数。证明下列条件等价：(i) p 不整除 \(\pi_{1}(G)\) 的挠子群的阶；
\end{enumerate}

\begin{enumerate}
\def\labelenumi{(\roman{enumi})}
\setcounter{enumi}{1}
\tightlist
\item
  对 G 的每个阶为 p 的元素 g，g 在 G 中的中心化子都是连通的；(iii) G 的每个同构于 \((\mathbf{Z}/p\mathbf{Z})^{2}\) 的子群都包含于一个极大环面。
\end{enumerate}

（为证明 (i) \(\Longrightarrow\) (ii)，利用习题 3 \(c\))；为证明 (iii) \(\Longrightarrow\) (i)，利用习题 6。）

\begin{enumerate}
\def\labelenumi{\arabic{enumi})}
\setcounter{enumi}{7}
\tightlist
\item
  取 \(\mathrm{G} = \mathbf{SO}(8, \mathbf{R}) / \{\pm I_8\}\)；用 \((\alpha_i)_{1 \leq i \leq 4}\) 表示 \(\mathrm{R}(\mathrm{G}, \mathrm{T})\) 的一个基，使得 \(\alpha_1, \alpha_3\) 与 \(\alpha_4\) 都不与 \(\alpha_2\) 正交（第 VI 章，图版 IV）。设 A 是 T 中由满足下列条件的 \(t \in T\) 组成的子群
\end{enumerate}

\[
\alpha_ {1} (t) = \alpha_ {3} (t) = \alpha_ {4} (t), \alpha_ {1} (t) ^ {2} = \alpha_ {2} (t) ^ {2} = 1.
\]

证明 A 同构于 \((\mathbf{Z}/2\mathbf{Z})^{2}\)，且它在 G 中的中心化子是一个有限的非交换群。

\begin{enumerate}
\def\labelenumi{\arabic{enumi})}
\setcounter{enumi}{8}
\tightlist
\item
  设 R 是一个不可约根系，\(R^{\vee}\) 是逆根系，B 是 R 的一个基，\(\alpha\) 是 R 的（相对于 B 的）最高根。置
\end{enumerate}

\[
\alpha = \sum_ {\beta \in \mathrm{B}} n _ {\beta} \beta \text { and } \alpha^ {\vee} = \sum_ {\beta \in \mathrm{B}} n _ {\beta} ^ {\vee} \beta^ {\vee};
\]

设 \(\nu (\mathrm{R}) = \sup_{\beta \in \mathrm{B}}n_{\beta}^{\vee}\)

\begin{enumerate}
\def\labelenumi{\alph{enumi})}
\item
  证明 N 中的区间 \(\mathbf{1},\nu(\mathrm{R})\mathbf{I}\) 是 1 与诸 \(n_{\beta}^{\vee}\)（\(\beta\in B\)）的并（置 \(\alpha_{0}=\alpha\)；设 \((\alpha_{1},\ldots,\alpha_{q})\) 是由 B 的互不相同的元素组成的一个序列，使得 \(\alpha_{i}\) 不与 \(\alpha_{i-1}\) 正交（对 \(i=1,\ldots,q\)），使得 \(n_{\alpha_{q}}^{\vee}=\nu(\mathrm{R})\)，且 q 对这些性质而言是极大的；证明 \(n_{\alpha_{i}}^{\vee}=i+1\)（对 \(i=1,\ldots,q\)））。
\item
  设 \(p\) 是一个素数。证明下列三条性质等价：
\end{enumerate}

\begin{enumerate}
\def\labelenumi{(\roman{enumi})}
\item
  \(p \leq \nu(\mathrm{R});\)
\item
  存在 \(\beta \in \mathrm{B}\) 使 \(p = n_{\beta}^{\vee}\)；
\item
  存在 \(\beta \in \mathrm{B}\) 使 \(p\) 整除 \(n_{\beta}^{\vee}\)。
\end{enumerate}

此时称 p 为 R 的一个挠素数。

\begin{enumerate}
\def\labelenumi{\alph{enumi})}
\setcounter{enumi}{2}
\tightlist
\item
  对不可约根系的每种型，给出 \(\nu(\mathrm{R})\) 的值以及挠素数。证明诸 \(n_{\beta}\) 组成的集合与诸 \(n_{\beta}^{\vee}\) 组成的集合除 \(G_{2}\) 型外是一致的。
\end{enumerate}

\(d\) ) 设 \(\mathbf{R}'\) 是 \(\mathbf{R}\) 的一个作为根系而言不可约的对称闭子集。证明 \(\nu (\mathrm{R}^{\prime})\leq \nu (\mathrm{R})\)

\begin{enumerate}
\def\labelenumi{\arabic{enumi})}
\setcounter{enumi}{9}
\tightlist
\item
  设 \(\mathbf{R}\) 是一个根系，\(p\) 是一个素数。证明下列性质等价：
\end{enumerate}

\begin{enumerate}
\def\labelenumi{(\roman{enumi})}
\item
  p 是 R 的某个不可约分支的挠素数（习题 9 b)）；\\
\item
  存在 R 的一个对称闭子集 \(R_{1}\)，它异于 R 且对这些性质而言是极大的，使得 \((\mathrm{Q}(\mathrm{R}^{\vee}):\mathrm{Q}(\mathrm{R}_{1}^{\vee}))=p(\mathrm{Q}(\mathrm{R}^{\vee}))\) 表示由 R 的逆根生成的 Z-模，而 \(\mathrm{Q}(\mathrm{R}_{1}^{\vee})\) 表示由诸 \(\alpha^{\vee}\)（\(\alpha\in R_{1}\)）生成的子模）；
\item
  存在 R 的一个对称闭子集 \(R_{1}\)，使 \(\mathrm{Q}(\mathrm{R}^{\vee})/\mathrm{Q}(\mathrm{R}_{1}^{\vee})\) 的 p-挠子模非零。
\end{enumerate}

（为证明 (i) \(\Longrightarrow\) (ii)，利用第 VI 章，§4 的习题 4；为证明 (iii) \(\Longrightarrow\) (i)，利用习题 9 d)。）

此时称 \(p\) 为 \(\mathbb{R}\) 的一个挠素数。

\begin{enumerate}
\def\labelenumi{\arabic{enumi})}
\setcounter{enumi}{10}
\tightlist
\item
  设 p 是一个素数。若存在 G 的一个具有极大秩的连通闭子群 H，使 \(\pi_{1}(H)\) 的 p-挠子群非零，则称 p 为 G 的一个挠素数。置 \(R = R(G, T)\)。
\end{enumerate}

\begin{enumerate}
\def\labelenumi{\alph{enumi})}
\item
  G 的挠素数是 R 的挠素数（习题 10）以及整除 \(\pi_1(\mathrm{G})\) 的挠群的阶的素数。
\item
  证明 G 的每个挠素数都整除 \(w/l!\)，其中 w 是 W 的阶，l 是半单群 D(G) 的秩（利用第 VI 章，§2，节 4，命题 7）。
\item
  设 \(t \in T\) 与 \(n \in N\) 满足 \(t^{n} \in \mathrm{C}(\mathrm{G})\)。设 \(R_{1}\) 是由 \(\alpha \in R\)（满足 \(\alpha(t) = 1\)）组成的集合，m 是 \(\mathrm{Q}(\mathrm{R}^{\vee}) / \mathrm{Q}(\mathrm{R}_{1}^{\vee})\) 的挠群的阶（参见~习题 10）。证明 m 的每个素因子都整除 n（化归到 G 单连通且概单的情形）。
\item
  假定 G 是单连通的。设 \(g \in G\) 与 \(n \in N\) 满足 \(g^{n}\) 属于 \(\mathrm{C}(\mathrm{G})\)，且 G 的任何挠素数都不整除 n。~证明 g 的中心化子的换位子群是单连通的。
\item
  设 g 是 G 的一个元素，p 是一个素数，满足 \(g^{p} = e\) 且 p 不是 G 的挠素数。证明中心化子 \(Z(g)\) 是连通的，且 p 不是 \(Z(g)\) 的挠素数。
\item
  假定 G 是单连通的，并设 p 是 G 的一个挠素数。证明存在 G 的一个阶为 p 的元素 g，使 \(\pi_{1}(\mathrm{Z}(g))\) 是 p 阶循环群（利用习题 10 以及第 VI 章，§4 的习题 4）。
\end{enumerate}

\begin{enumerate}
\def\labelenumi{\arabic{enumi})}
\setcounter{enumi}{11}
\tightlist
\item
  设 \(p\) 是一个素数。证明下列条件等价：
\end{enumerate}

\begin{enumerate}
\def\labelenumi{(\roman{enumi})}
\item
  \(p\) 不是 \(\mathbf{G}\) 的挠素数；
\item
  对每个子群 \(\mathbf{F}\)（属于 \(\mathbf{G}\)，同构于 \((\mathbf{Z} / p\mathbf{Z})^n\)（\(n\in \mathbf{N}\)）），\(\mathbf{F}\) 在 \(\mathbf{G}\) 中的中心化子都是连通的；
\end{enumerate}

(ii') 对 G 的每个同构于 \((\mathbf{Z}/p\mathbf{Z})^{2}\) 的子群 F，F 在 G 中的中心化子都是连通的；

\begin{enumerate}
\def\labelenumi{(\roman{enumi})}
\setcounter{enumi}{2}
\tightlist
\item
  G 的每个同构于 \((\mathbf{Z}/p\mathbf{Z})^{n}\)（n 为某个整数）的子群都包含于一个极大环面；
\end{enumerate}

(iii') G 的每个同构于 \((\mathbf{Z}/p\mathbf{Z})^{3}\) 的子群都包含于一个极大环面。

（为证明 (i) \(\Longrightarrow\) (ii)，利用习题 11 e)；为证明 (iii) \(\Longrightarrow\) (i)，利用习题 11 f) 以及习题 7。）

\exercisesection{\S~6}

\begin{enumerate}
\def\labelenumi{\arabic{enumi})}
\tightlist
\item
  设 R 是实向量空间 V 中的一个既约根系。给空间 V 赋以一个在 W(R) 下不变的内积，并给空间 S(V) 赋以相应的内积（《拓扑向量空间》第 V 章，§3，节 3，公式 (13) 与 (14)）；于是，对 V 中的 \(x_{1},\ldots,x_{n},y_{1},\ldots,y_{n}\)，
\end{enumerate}

\[
(x _ {1} \dots x _ {n} | y _ {1} \dots y _ {n}) = \sum_ {\sigma \in \mathfrak {S} _ {n}} (x _ {1} | y _ {\sigma (1)}) \dots (x _ {n} | y _ {\sigma (n)}).
\]

选取 R 的一个房 C；置 \(N = \text{Card}(R_{+})\)，\(\rho = \frac{1}{2} \sum_{\alpha \in R_{+}} \alpha\) 与 \(w(R) = \text{Card}(W(R))\)。设 P 是元素 \(\prod_{\alpha \in R_{+}} \alpha\)（\(\mathbf{S}^{\mathrm{N}}(\mathrm{V})\) 中的）。

\(a)\) 证明 \(\mathrm{P} = \frac{1}{\mathrm{N!}}\sum_{w\in \mathrm{W(R)}}\varepsilon (w)(w\rho)^{\mathrm{N}}\)（参见~第 VI 章，§3，节 3，命题 2）。

\begin{enumerate}
\def\labelenumi{\alph{enumi})}
\setcounter{enumi}{1}
\item
  推出等式 \((\mathrm{P}|\mathrm{P}) = w(\mathrm{R})\prod_{\alpha \in \mathrm{R}_{+}}(\rho |\alpha)\)。
\item
  证明 \((\mathrm{P}|\mathrm{P}) = 2^{-\mathrm{N}}w(\mathrm{R})\prod_{i = 1}^{l}m_i!\prod_{\alpha \in \mathrm{R}_+}(\alpha |\alpha) = 2^{-\mathrm{N}}\prod_{i = 1}^{l}(m_i + 1)!\prod_{\alpha \in \mathrm{R}_+}(\alpha |\alpha)\)，其中 \(m_{1},\ldots ,m_{l}\) 是 \(\mathrm{W(R)}\) 的指数（利用第 VIII 章，§9 的习题 3）。
\end{enumerate}

¶ 2) 设 V 是 R 上的一个有限维实希尔伯特空间。给 V 上的多项式空间 \(\mathbf{S}(\mathrm{V}^{*})\) 赋以《拓扑向量空间》第 V 章，§3，节 3，公式 (13) 与 (14) 所定义的内积（参见~习题 1）。用 \(\gamma\) 表示 V 上的典则高斯测度（《积分》第 IX 章，§6，节 4 至 6）；于是，若 \((x_{i})_{1\leq i\leq n}\) 是 \(V^{*}\) 的一个标准正交基，则有 \(d\gamma(x)=(2\pi)^{-n/2}e^{-(x|x)/2}dx_{1}\ldots dx_{n}\)。

\begin{enumerate}
\def\labelenumi{\alph{enumi})}
\tightlist
\item
  设 \(q\) 是 \(\mathbf{S}^2(\mathrm{V})\) 中定义 \(\mathrm{V}^*\) 上内积的元素，并设 \(\Delta: \mathbf{S}(\mathrm{V}^*) \to \mathbf{S}(\mathrm{V}^*)\) 是由与 \(q\) 作内积所给出的算子（《代数》第 III 章，§11，节 9）。证明对每个标准正交基 \((x_i)_{1 \leq i \leq n}\)（\(\mathrm{V}^*\) 的）以及每个多项式 \(\mathrm{P} \in \mathbf{S}(\mathrm{V}^*)\)，有 \(\Delta(\mathrm{P}) = \frac{1}{2} \sum_{i=1}^{n} \frac{\partial^2\mathrm{P}}{\partial x_i^2}\)。b) 对 \(\mathrm{P} \in \mathbf{S}(\mathrm{V}^*)\)，置 \(\mathrm{P}^* = \mathrm{P} * \gamma\)，使得 \(\mathrm{P}^*(x) = \int_{\mathrm{V}} \mathrm{P}(x - y) d\gamma(y)\)（对 \(x \in \mathrm{V}\)）。证明等式 \(\mathrm{P}^* = \sum_{n=1}^{\infty} \frac{\Delta^n(\mathrm{P})}{n!} = e^{\Delta}(\mathrm{P})\)。
\end{enumerate}

（化归为对函数 \(x \mapsto e^{i(x|u)}\)（\(u \in \mathrm{V}\)）证明类似的公式。）

\begin{enumerate}
\def\labelenumi{\alph{enumi})}
\setcounter{enumi}{2}
\tightlist
\item
  证明公式 \(\int\overline{\mathrm{P}^{*}(ix)}\mathrm{Q}^{*}(ix)d\gamma(x)=(\mathrm{P}|\mathrm{Q})\)（对 \(\mathbf{S}(\mathrm{V}^{*})\) 中的 P 与 Q，后者等同于 \(\mathbf{S}_{\mathbf{C}}((\mathrm{V}\otimes\mathbf{C})^{*})\) 的一个子空间）。
\end{enumerate}

\(d)\) 若 \(\mathrm{P}\) 是 \(\mathrm{V}\) 上的一个齐次多项式且 \(\Delta (\mathrm{P}) = 0\)，则 \(\int_{\mathrm{V}}\mathrm{P}(x)^2 d\gamma (x) = (\mathrm{P}|\mathrm{P})\)。

\begin{enumerate}
\def\labelenumi{\alph{enumi})}
\setcounter{enumi}{4}
\tightlist
\item
  设 W 是 V 的一个由反射生成的有限自同构群；用 H 表示 W 中反射组成的集合。对 \(h \in H\)，设 \(e_{h} \in V\) 与 \(f_{h} \in V^{*}\) 满足 \(h(x) = x + f_{h}(x)e_{h}\)。证明多项式 \(P = \prod_{h \in H} f_{h}\) 满足 \(\Delta(P) = 0\)（利用第 V 章，§5，节 4，命题 5）。
\end{enumerate}

\begin{enumerate}
\def\labelenumi{\arabic{enumi})}
\setcounter{enumi}{2}
\tightlist
\item
  设 \(\mu\) 是加法群 \(\mathfrak{g}\) 上的一个 Haar 测度。
\end{enumerate}

\begin{enumerate}
\def\labelenumi{\alph{enumi})}
\item
  G 上存在唯一的具有下列性质的 Haar 测度 \(\mu_{G}\)：若 \(\omega_{G}\) 与 \(\omega_{g}\) 分别是 G 与 g 上次数为 n 的不变微分形式，满足 \(\omega_{\mathrm{G}}(e)=\omega_{\mathfrak{g}}(0)\) 且 \(|\omega_{g}|=\mu\)，则 \(|\omega_{G}|=\mu_{G}\)。映射 \(\mu\mapsto\mu_{G}\) 是从 g 上 Haar 测度的集合到 G 上相应的集合的一个双射。
\item
  选取一个在 g 下不变的内积 ( \textbar{} )，以及 t 上的一个 Haar 测度 \(\tau\)，使得 \(\mu\)（相应地 \(\tau\)）当 g（相应地 t）借助一个标准正交基与 \(R^{n}\)（相应地与 \(R^{r}\)）等同起来时对应于 Lebesgue 测度。证明公式
\end{enumerate}

\[
\int_ {t} \pi_ {\mathfrak {g}} (x) e ^ {- (x | x) / 2} d \tau (x) = (2 \pi) ^ {n / 2} \frac {\tau_ {\mathrm{T}} (\mathrm{T})}{\mu_ {\mathrm{G}} (\mathrm{G})} w (\mathrm{G}).
\]

\begin{enumerate}
\def\labelenumi{\alph{enumi})}
\setcounter{enumi}{2}
\tightlist
\item
  把 X(T) 与希尔伯特空间 \(t^{*}\) 的一个子集借助映射 \((2\pi i)^{-1}\delta\)（§4，节 2）等同起来，并置 \(\mathrm{P}(x)=\prod_{\alpha\in\mathrm{R}_{+}}\langle\alpha,x\rangle\)（对 \(x\in t\)）。用习题 1 与 2 的记号，证明
\end{enumerate}

\[
\frac {\tau_ {\mathrm{T}} (\mathrm{T})}{\mu_ {\mathrm{G}} (\mathrm{G})} = (2 \pi) ^ {\mathrm{N}} w (\mathrm{G}) ^ {- 1} (\mathrm{P} | \mathrm{P}) = \pi^ {\mathrm{N}} \prod_ {\alpha \in \mathrm{R} _ {+}} (\alpha | \alpha) \prod_ {i = 1} ^ {l} m _ {i}!
\]

\(d\) ) 用 §3 习题 9 的记号，设 \(\mathfrak{g}_{\mathbf{Z}}\) 是 \(\mathfrak{g}\) 中由 \((2\pi)^{-1}\Gamma (\mathrm{T})\) 与诸元素 \(u_{\alpha},v_{\alpha}\)（\(\alpha \in \mathbb{R}\)）生成的 Z-李子代数。证明等式

\[
\mu_ {\mathrm{G}} (\mathrm{G}) = \mu (\mathfrak {g} / \mathfrak {g} _ {\mathbf {Z}}) \frac {2 ^ {r} \pi^ {\mathrm{N} + r}}{\prod_ {i} m _ {i} !}.
\]

\begin{enumerate}
\def\labelenumi{\alph{enumi})}
\setcounter{enumi}{4}
\tightlist
\item
  假定 G 是单连通的。证明 l = r，且
\end{enumerate}

\[
\mu_ {\mathrm{G}} (\mathrm{G}) = 2 ^ {l / 2} \pi^ {- \mathrm{N}} f ^ {1 / 2} \prod_ {\alpha \in \mathrm{R} _ {+}} (\alpha | \alpha) ^ {- 1} \prod_ {i = 1} ^ {l} (\alpha_ {i} | \alpha_ {i}) ^ {- 1 / 2} \prod_ {i} (m _ {i}!) ^ {- 1},
\]

其中 \(\{\alpha_1, \ldots, \alpha_l\}\) 是 \(R\) 的一个基，\(f\) 是 \(R\) 的连通指数。

\(f\) ) 再假定 \(\mathbb{R}\) 不可约且其所有根都等长；取 \(\mathfrak{g}\) 上的内积为 Killing 型的相反数。证明 \(\mu_{\mathrm{G}}(\mathrm{G}) = (2\pi)^{\mathrm{N} + r}(2h)^{n / 2}f^{1 / 2}\prod_{i}(m_i!)^{-1}\)。

\begin{enumerate}
\def\labelenumi{\arabic{enumi})}
\setcounter{enumi}{3}
\tightlist
\item
  设 \(\mathrm{X}\) 是一个 \(\mathrm{C}^{\infty}\) 类微分流形。在本习题及以下诸习题中，我们将用 \(\mathrm{H}(\mathrm{X})\) 简单地表示分次 \(\mathbf{R}\)-空间 \(\mathrm{H}(\Omega (\mathrm{X}))\)。
\end{enumerate}

\begin{enumerate}
\def\labelenumi{\alph{enumi})}
\item
  证明 \(\Omega(\mathrm{X})\) 是一个结合且反交换的分次微分代数（《代数学》（Algèbre）第 X 章第~183 页，习题 18）。由此推出 \(\mathrm{H}(\mathrm{X})\) 具有一个自然的结合且反交换的分次代数结构。
\item
  若 X 连通且维数为 p，则 \(\mathrm{H}^{i}(\mathrm{X}) = 0\)（对 i \textgreater{} p），且 \(\dim_{\mathbf{R}} \mathrm{H}^{0}(\mathrm{X}) = 1\)。此外，若 X 是紧的，则空间 \(\mathrm{H}^{p}(\mathrm{X})\) 的维数为 1（利用《微分流形与解析流形，结果》，11.2.4）。
\item
  设 Y 是另一个 \(C^{\infty}\) 类流形，\(f: X \to Y\) 是一个 \(C^{\infty}\) 类映射。映射 \(f^{*}: \Omega(Y) \to \Omega(X)\) 是一个复形态射（《微分流形与解析流形，结果》，8.3.5）；证明 \(\mathrm{H}(f^{*}) : \mathrm{H}(Y) \to \mathrm{H}(X)\) 是一个代数态射。
\item
  假定给定了 G 在 X 上的一个 \(C^{\infty}\) 类作用律。证明由不变形式组成的子复形 \(\Omega(\mathrm{X})^{\mathrm{G}}\) 是 \(\Omega(\mathrm{X})\) 的一个子代数。由此推出定理 2 中定义的映射 \(\mathrm{H}(i):\mathrm{H}(\Omega(\mathrm{X})^{\mathrm{G}})\to\mathrm{H}(\mathrm{X})\) 是分次代数的一个同构。
\item
  证明 \((\operatorname{Alt}(\mathfrak{g}))^{\mathrm{G}}\) 是 \(\operatorname{Alt}(\mathfrak{g})\) 的一个子代数，且分次代数 \(\operatorname{H}(\operatorname{G})\) 与之同构。
\end{enumerate}

\begin{enumerate}
\def\labelenumi{\arabic{enumi})}
\setcounter{enumi}{4}
\tightlist
\item
  用 H(G) 表示分次 R-代数 H(Ω(G))（参见~习题 4）。
\end{enumerate}

\(a\) ) 空间 \(\mathrm{H}^p (\mathbf{G})\) 当 \(p > n\) 时为零，当 \(p = n\) 与 \(p = 0\) 时维数为 1。空间 \(\mathrm{H}^1 (\mathrm{G})\) 可典则地与 \(\mathfrak{c}^*\) 等同，其中 \(\mathfrak{c} = \mathrm{L}(\mathrm{C}(\mathrm{G}))\)。

\begin{enumerate}
\def\labelenumi{\alph{enumi})}
\setcounter{enumi}{1}
\item
  以下假定 G 是半单的。证明 \(\mathrm{H}^{1}(\mathrm{G})=\mathrm{H}^{2}(\mathrm{G})=0\)（参见~第 I 章，§6，习题 1）。
\item
  用 \(\mathrm{B}(\mathfrak{g})\) 表示 \(\mathfrak{g}\) 上 G-不变的对称双线性型组成的空间；对 \(b \in \mathrm{B}(\mathfrak{g})\) 与 \(x, y, z\)（属于 \(\mathfrak{g}\)），置 \(\tilde{b}(x, y, z) = b([x, y], z)\)。证明映射 \(b \mapsto \tilde{b}\) 定义了从 \(\mathrm{B}(\mathfrak{g})\) 到 \(\mathrm{H}^{3}(\mathrm{G})\) 的一个同构（设 \(\omega\in\mathrm{H}^{3}(\mathrm{G})\)；证明对所有 \(x\in g\)，存在 g 上唯一的线性形式 \(f(x)\) 使 \(df(x)=i(x)\omega\)，并考察形式 \((x,y)\mapsto-\langle y,f(x)\rangle\)）。
\end{enumerate}

\(d\)) 证明 \(\mathbf{R}\)-向量空间 \(\mathrm{H}^3 (\mathrm{G})\) 的维数等于 \(\mathfrak{g}\) 的单理想的个数。

\begin{enumerate}
\def\labelenumi{\arabic{enumi})}
\setcounter{enumi}{5}
\tightlist
\item
  置 \(b_{i}(\mathrm{G}) = \dim_{\mathbf{R}} \mathrm{H}^{i}(\mathrm{G})\)（对 \(i \geq 0\)），且若 \(\mathrm{X}\) 表示一个不定元，置 \(\mathrm{P}_{\mathrm{G}}(\mathrm{X}) = \sum_{i \geq 0} b_{i}(\mathrm{G}) \mathrm{X}^{i}\)。
\end{enumerate}

\begin{enumerate}
\def\labelenumi{\alph{enumi})}
\item
  证明 \(b_{i}(\mathrm{G}) = \int_{\mathrm{G}} \mathrm{Tr} \wedge^{i} (\mathrm{Ad} g) dg\) 与 \(\mathrm{P}_{\mathrm{G}}(\mathrm{X}) = \int_{\mathrm{G}} \det(1 + \mathrm{X}. \mathrm{Ad} g) dg\)（参见~附录 II，引理 1）。
\item
  推出等式 \(\sum_{i} b_{i}(\mathrm{G}) = 2^{r}\)，以及 \(\sum_{i} (-1)^{i} b_{i}(\mathrm{G}) = 0\)（当 \(\dim \mathrm{G} > 0\) 时）（利用 H. Weyl 公式，或 §2 的习题 8）。
\item
  取 \(\mathrm{G} = \mathbf{U}(n, \mathbf{C})\)。证明 \(\mathrm{P}_{\mathrm{G}}(\mathrm{X})\) 是 \((\mathrm{X}_1 \ldots \mathrm{X}_n)^{2n-2}\) 在多项式 \(\frac{1}{n!}(1 + \mathrm{X})^n \prod_{\substack{1 \leq i, j \leq n \\ i \neq j}} (\mathrm{XX}_i + \mathrm{X}_j)(\mathrm{X}_i - \mathrm{X}_j)\)（系数在 \(\mathbf{Z}[\mathrm{X}]\) 中）中的系数。
\end{enumerate}

\begin{enumerate}
\def\labelenumi{\arabic{enumi})}
\setcounter{enumi}{6}
\tightlist
\item
  设 \(\mathbf{K}\) 是一个连通紧李群。
\end{enumerate}

\begin{enumerate}
\def\labelenumi{\alph{enumi})}
\item
  设 \(f: K \to G\) 是一个核为有限的满同态。证明同态 \(\mathrm{H}(f^{*}) : \mathrm{H}(\mathrm{G}) \to \mathrm{H}(\mathrm{K})\) 是一个同构。
\item
  证明代数 \(\mathrm{H}(\mathrm{G}\times\mathrm{K})\) 可典则地与斜张量积 \(\mathrm{H}(\mathrm{G})^{g}\otimes\mathrm{H}(\mathrm{K})\) 等同。
\item
  由 a) 与 b) 推出，作为代数，\(\mathrm{H}(\mathrm{G})\) 同构于 \(\mathrm{H}(\mathrm{C}(\mathrm{G})_{0})^{g}\otimes\mathrm{H}(\mathrm{D}(\mathrm{G}))\)。证明代数 \(\mathrm{H}(\mathrm{C}(\mathrm{G})_{0})\) 同构于 \(\wedge(\mathfrak{c}^{*})\)，其中 \(\mathfrak{c}=\mathrm{L}(\mathrm{C}(\mathrm{G}))\)。
\end{enumerate}

¶ 8) 设 k 是一个特征为零的域，E 是 k 上的一个分次左双代数（《代数》第 III 章，§11，节 4，定义 3）。假定 E 是反交换且反余交换的，且当 m 充分大时 \(E_{m} = 0\)。用 P 表示 E 的本原元组成的集合（参见~第 II 章，§1）。

\begin{enumerate}
\def\labelenumi{\alph{enumi})}
\item
  证明 P 的每个齐次元素都是奇次数的（对大的 m 展开 \(c(x^{m})\)）。由此推出存在分次左双代数的一个典则态射 \(\varphi : \Lambda(\mathrm{P}) \to \mathrm{E}\)（\(\Lambda(\mathrm{P})\) 的分次左双代数结构即《代数》第 III 章，§11，习题 6 中所定义的那一个）。
\item
  证明 \(\varphi\) 是一个同构（把第 II 章，§1，节 6 定理 1 的证明加以改造）。
\end{enumerate}

\begin{enumerate}
\def\labelenumi{\arabic{enumi})}
\setcounter{enumi}{8}
\tightlist
\item
  用 \(m: \mathrm{G} \times \mathrm{G} \to \mathrm{G}\) 表示满足 \(m(g, h) = gh\)（对 \(g, h\)，属于 \(\mathrm{G}\)）的映射。把 \(\mathrm{H}(\mathrm{G} \times \mathrm{G})\) 与 \(\mathrm{H}(\mathrm{G})^g \otimes \mathrm{H}(\mathrm{G})\) 等同起来（习题 7 b)），于是 \(m^*\) 定义了一个代数同态 \(c: \mathrm{H}(\mathrm{G}) \to \mathrm{H}(\mathrm{G})^g \otimes \mathrm{H}(\mathrm{G})\)。
\end{enumerate}

\begin{enumerate}
\def\labelenumi{\alph{enumi})}
\item
  证明 \((\mathrm{H}(\mathrm{G}), c)\) 是一个反交换且反余交换的分次左双代数（注意映射 \(g \mapsto g^{-1}\) 在 \(\mathrm{H}^{p}(\mathrm{G})\) 上诱导乘以 \((-1)^{p}\)）。
\item
  设 \(\mathrm{P}(\mathrm{G})\) 是 \(\mathrm{H}(\mathrm{G})\) 中由本原元素组成的分次子空间；由习题 8 推出存在分次双代数的一个同构 \(\Lambda(\mathrm{P}(\mathrm{G})) \to \mathrm{H}(\mathrm{G})\)。
\item
  证明 \(\dim_{\mathbf{R}}\mathrm{P}(\mathrm{G})=r\)（利用习题 6 b)）。
\end{enumerate}

\(d\) ) 推出多项式 \(\sum_{i\geq 0}b_i(\mathrm{G})\mathrm{X}^i\) 具有如下形状

\[
(1 + \mathrm{X}) ^ {c} \prod_ {i = 1} ^ {l} (1 + \mathrm{X} ^ {2 k _ {i} + 1}),
\]

其中 c 是 \(\mathrm{C}(\mathrm{G})\) 的维数，l 是 \(\mathrm{D}(\mathrm{G})\) 的秩，诸 \(k_{i}\) 是 \(\geq 1\) 的整数；我们有 \(k_{1} + \cdots + k_{l} = \frac{1}{2}\mathrm{Card}\mathrm{R}(\mathrm{G},\mathrm{T})\)\(^{12}\)。

\footnotetext[12]{整数 \(k_{i}\) 事实上是 \(\mathrm{R}(\mathrm{G},\mathrm{T})\) 的指数。参见 J. LERAY, Sur l'homologie des groupes de Lie, des espaces homogènes et des espaces fibrés principaux, Colloque de Topologie de Bruxelles (1950), pp. 101-115.}

\begin{enumerate}
\def\labelenumi{\arabic{enumi})}
\setcounter{enumi}{9}
\tightlist
\item
  设 \(\mathrm{H}\) 是 \(\mathrm{G}\) 的一个闭子群；用 \(dh\) 表示 \(\mathrm{H}\) 上总质量为 1 的 Haar 测度。置 \(\chi(\mathrm{G} / \mathrm{H}) = \sum_{i > 0} (-1)^i \dim_{\mathbf{R}} \mathrm{H}^i (\mathrm{G} / \mathrm{H})\)。
\end{enumerate}

\begin{enumerate}
\def\labelenumi{\alph{enumi})}
\item
  证明等式 \(\chi(\mathrm{G}/\mathrm{H})=\overline{\int_{\mathrm{H}}}\det(1-\mathrm{Ad}_{\mathfrak{g}/\mathfrak{h}}h)\,dh\)（利用《代数学》（Algèbre）第 X 章第~41 页，命题 11 以及附录 II，节 2，引理 1）。
\item
  设 \(\pi_{0}(\mathrm{H})\) 是 H 的连通分支的个数。证明
\end{enumerate}

\(\chi (\mathrm{G / H}) = 0\)，若 \(\mathrm{H_0}\) 不是极大秩

\(\chi (\mathrm{G / H}) = w(\mathrm{G}) / w(\mathrm{H}_0)\pi_0(\mathrm{H})\)，若 \(\mathrm{H_0}\) 是极大秩。

\begin{enumerate}
\def\labelenumi{\arabic{enumi})}
\setcounter{enumi}{10}
\tightlist
\item
  设 \(u\) 是 \(G\) 的一个阶为 2 的自同构；用 \(K\) 表示 \(u\) 的不动点集合的单位元连通分支，\(\mathfrak{k}\) 表示它的李代数，\(X\) 表示对称空间 \(G / K\)（§1，习题 8）。
\end{enumerate}

\begin{enumerate}
\def\labelenumi{\alph{enumi})}
\item
  证明 X 上每个 G-不变微分形式 \(\omega\) 都满足 \(d\omega = 0\)（注意 \(u\) 在 \(\mathrm{Alt}^p (\mathfrak{g} / \mathfrak{k})\) 上诱导乘以 \((-1)^p)\)）。
\item
  由此推出，从分次代数 \(\mathrm{H}(\mathrm{G}/\mathrm{K})\) 到由 K-不变元素组成的 \(\operatorname{Alt}(\mathfrak{g}/\mathfrak{k})\) 的分次子代数有一个同构。
\item
  置 \(b_{i}(\mathrm{G / K}) = \dim_{\mathbf{R}}\mathrm{H}^{i}(\mathrm{G / K})\)（对 \(i\geq 0\)）；对 \(k\in \mathrm{K}\)，用 \(\mathrm{Ad}^{-}k\) 表示 \(\mathrm{Ad}k\) 在 \(\mathrm{L}(u)\) 的特征值为 \(-1\) 的特征子空间上的限制。设 \(dk\) 是 \(\mathbf{K}\) 上总质量为 1 的 Haar 测度。证明公式 \(b_{i}(\mathrm{G / K}) = \int_{\mathrm{K}}\mathrm{Tr}\wedge^{i}(\mathrm{Ad}^{-}k)dk\) 与 \(\sum_{i > 0}b_{i}(\mathrm{G / K})\mathrm{X}^{i} = \int_{\mathrm{K}}\det (1 + \mathrm{XAd}^{-}k)dk\)。
\item
  证明若此外 K 还是极大秩的，则代数 \(\mathrm{H}(\mathrm{G}/\mathrm{K})\) 的奇数部分为零（注意 \(u = \operatorname{Int} k\)，其中 \(k \in K\)）。
\item
  计算分次代数 \(\mathrm{H}(\mathbf{S}_{n})\)。由此推出 \(S_{n}\) 具有一个李群结构（与其流形结构相容）当且仅当 n 等于 1 或 3。
\end{enumerate}

\begin{enumerate}
\def\labelenumi{\arabic{enumi})}
\setcounter{enumi}{11}
\tightlist
\item
  设 H 是 G 的一个连通闭子群，h 是它的李代数。
\end{enumerate}

\begin{enumerate}
\def\labelenumi{\alph{enumi})}
\item
  设 \(\alpha\) 是 \(\mathfrak{h}^*\) 的一个在 H 下不变的元素；证明存在元素 \(\bar{\alpha}\)（属于 \(\mathfrak{g}^*\)，在 H 下不变，且它到 \(\mathfrak{h}\) 上的限制等于 \(\alpha\)）。证明元素 \(d\bar{\alpha} \in \mathrm{Alt}^2(\mathfrak{g})\) 被 \(i(\eta)\) 与 \(\theta(\eta)\)（对所有 \(\eta \in \mathfrak{h}\)）零化，且它在 \(\mathrm{H}^2(\mathrm{G}/\mathrm{H})\) 中的类不依赖于 \(\bar{\alpha}\) 的选取。这就定义了一个同态 \(\varphi: \mathrm{H}^1(\mathrm{H}) \to \mathrm{H}^2(\mathrm{G}/\mathrm{H})\)。
\item
  以下假定 \(\mathbf{G}\) 是半单的。证明 \(\varphi\) 是一个同构。
\item
  证明 \(\mathrm{H}\) 是半单的当且仅当 \(\mathrm{H}^2 (\mathrm{G / H}) = 0\)。
\end{enumerate}

\(d\) ) 不假定 \(\mathrm{H}\) 连通，定义一个同构 \((\mathfrak{h}^{*})^{\mathrm{H}}\to\) \(\mathrm{H}^2 (\mathrm{G / H})\)（把 \(b\) ) 应用于 \(\mathrm{H}_0\)）。

\begin{enumerate}
\def\labelenumi{\arabic{enumi})}
\setcounter{enumi}{12}
\tightlist
\item
  设 \(\mathrm{H}\) 是 \(\mathrm{G}\) 的一个闭子群；置 \(\mathrm{X} = \mathrm{G} / \mathrm{H}\)，\(n = \dim \mathrm{X}\)。证明下列条件等价：
\end{enumerate}

\begin{enumerate}
\def\labelenumi{(\roman{enumi})}
\item
  在 X 上存在一个 2-形式 \(\omega\)，使得 \(d\omega = 0\)，且使得交错形式 \(\omega_{x}\)（\(\mathrm{T}_x(\mathrm{X})\) 上的）对所有 \(x\in \mathrm{X}\) 都是分离的；
\item
  \(n\) 是偶数，且存在元素 \(\omega\)（属于 \(\mathrm{H}^2 (\mathrm{G / H})\)）使 \(\omega^{n / 2}\neq 0\)；(iii) \(\mathrm{H}\) 是 \(\mathbf{G}\) 中一个环面的中心化子；
\item
  H 是极大秩的，且在 X 上存在一个 G-不变的复结构；
\item
  在 X 上存在一个复结构 \(j\) 与一个 2-形式 \(\omega\)，使得 \(d\omega = 0\)，\(\omega_{x}(ju,jv) = \omega_{x}(u,v)\)，且 \(\omega_{x}(u,ju) > 0\)（对所有 X 中的 \(x\) 以及所有非零的 \(u,v\)（属于 \(\mathrm{T}_x(\mathrm{X})\)））（* 换言之，即 \(\mathrm{X}_{*}\) 上的一个 Kähler 结构*）；
\item
  在 X 上存在满足 (v) 中条件且在 G 下不变的复结构 j 与 2-形式 \(\omega\)（* 即 \(X_{*}\) 上的一个 G-不变 Kähler 结构*）。
\end{enumerate}

（为证明 (ii) \(\Longrightarrow\) (iii)，置 \(S = C(H)_{0}\) 与 \(Z = Z_{G}(S)\)；利用习题 12 d)，证明典则映射 \(H^{2}(G/Z) \to H^{2}(G/H)\) 是满射，并由此推出 Z = H。(iii) 与 (iv) 的等价性由 §4 的习题 8 得出。为证明 (iii) \(\Longrightarrow\) (vi)，在 g/L(H) 上构造一个 H-不变的分离正埃尔米特形式并考察其虚部；利用习题 11 a)。）

\exercisesection{\S~7}

\begin{enumerate}
\def\labelenumi{\arabic{enumi})}
\tightlist
\item
  取 G 为群 \(\mathbf{U}(n,\mathbf{C})\)，T 为由对角矩阵组成的子群；对 T 中的 \(t = \text{diag}(t_1, \ldots, t_n)\) 与 \(1 \leq i \leq n\)，置 \(\varepsilon_i(t) = t_i\)。用 \(\sigma\) 表示 G 在 \(C^n\) 上的恒等表示。
\end{enumerate}

\begin{enumerate}
\def\labelenumi{\alph{enumi})}
\item
  群 X(T) 有基 \(\varepsilon_{1},\ldots,\varepsilon_{n}\)；证明 \(\mathbf{Z}[\mathrm{X}(\mathrm{T})]^{\mathrm{W}}\) 的每个元素都具有形状 \(e^{k(\varepsilon_{1}+\cdots+\varepsilon_{n})}\mathrm{P}(e^{\varepsilon_{1}},\ldots,e^{\varepsilon_{n}})\)，其中 k 是一个相对整数，P 是 n 个变量的、系数为整数的对称多项式。
\item
  证明表示 \(\Lambda^r\sigma\)（\(1 \leq r \leq n\)）都是不可约的。
\item
  证明同态 \(u: \mathbf{Z}[X_1, X_2, \ldots, X_n][X_n^{-1}] \to \mathrm{R}(\mathrm{G})\)（满足 \(u(X_i) = [\Lambda^i \sigma]\)）是一个同构。
\end{enumerate}

\begin{enumerate}
\def\labelenumi{\arabic{enumi})}
\setcounter{enumi}{1}
\tightlist
\item
  设 \(\mathrm{G} = \mathbf{SO}(2l + 1, \mathbf{R})\)，其中 \(l \geq 1\)；我们使用 §4 习题 2 的记号。用 \(\sigma\) 表示 \(\mathrm{G}\) 在 \(\mathbf{C}^{2l + 1}\) 上由恒等表示通过标量扩张得到的表示。
\end{enumerate}

\begin{enumerate}
\def\labelenumi{\alph{enumi})}
\item
  Z-模 X(T) 以 \(\varpi_{1},\ldots,\varpi_{l-1},2\varpi_{l}\) 为基，同样也以 \(\varepsilon_{1},\ldots,\varepsilon_{l}\) 为基。
\item
  用 \(\eta_{i}\) 表示元素 \(e^{\varepsilon_{i}} + e^{-\varepsilon_{i}}\)（属于 Z{[}X(T){]}）。证明 \(\mathbf{Z}[X(T)]^{W}\) 的每个元素都可写成 \(\mathrm{P}(\eta_{1}, \ldots, \eta_{l})\)，其中 P 是 l 个变量的、系数为整数的对称多项式。
\item
  证明表示 \(\Lambda^r\sigma (r\leq 2l + 1)\) 都是不可约的。对 \(r\leq l\)，证明等式
\end{enumerate}

\[
\begin{array}{c} \operatorname{Ch} (\Lambda^ {r}   \sigma) = s _ {r} (\eta_ {1}, \ldots , \eta_ {l}) + (l - r + 2) s _ {r - 2} (\eta_ {1}, \ldots , \eta_ {l}) + \dots + \\ \qquad \qquad + \binom {l - r + 2 k} {k} s _ {r - 2 k} (\eta_ {1}, \ldots , \eta_ {l}) + \dots , \end{array}
\]

其中诸 \(s_{k}\) 是 l 个变量的初等对称多项式。

\(d\) ) 证明同态 \(u:\mathbf{Z}[\mathrm{X}_1,\dots ,\mathrm{X}_l]\to \mathrm{R}(\mathrm{G})\)（满足 \(u(\mathrm{X}_i) = [\Lambda^i\sigma ]\)）是一个同构。

\begin{enumerate}
\def\labelenumi{\arabic{enumi})}
\setcounter{enumi}{2}
\tightlist
\item
  设 \(\mathrm{G} = \mathbf{SO}(2l, \mathbf{R})\)，其中 \(l \geq 2\)；我们使用 §4 习题 2 的记号。用 \(\sigma\) 表示 G 在 \(C^{2l}\) 上由恒等表示通过标量扩张得到的表示。
\end{enumerate}

\begin{enumerate}
\def\labelenumi{\alph{enumi})}
\item
  \(\mathbf{Z}\)-模 \(\mathrm{X(T)}\) 有基 \(\varepsilon_1, \ldots, \varepsilon_l\)。
\item
  用 \(\eta_i\) 表示元素 \(e^{\varepsilon_i} + e^{-\varepsilon_i}\)（属于 \(\mathbf{Z}[\mathrm{X(T)}]\)），用 \(\delta\) 表示元素 \(\prod_{i=1}^{l}(e^{\varepsilon_i} - e^{-\varepsilon_i})\)。证明 \(\mathbf{Z}[\mathrm{X(T)}]^{\mathrm{W}}\) 的每个元素都可写成 \(\mathrm{P}(\eta_1, \ldots, \eta_l) + \mathrm{Q}(\eta_1, \ldots, \eta_l)\delta\)，其中 \(\mathrm{P}\) 与 \(\mathrm{Q}\) 是 \(l\) 个变量的、系数在 \(\frac{1}{2}\mathbf{Z}\) 中的对称多项式。
\item
  证明表示 \(\Lambda^{r}\sigma\) 当 \(r\neq l\) 时都是不可约的；表示 \(\Lambda^{l}\sigma\) 是两个子表示 \(\tau^{+}\) 与 \(\tau^{-}\) 的直和，其最高权分别为 \(2\varpi_{l}\) 与 \(2\varpi_{l-1}\)（见第 VIII 章，§13，习题 10）。
\end{enumerate}

\(d\) ) 证明当 \(r \leq l\) 时，元素 \(\operatorname{Ch}(\Lambda^r \sigma)\) 由习题 2 \(c\)) 中的公式给出，且

\[
\mathrm{Ch} (\tau^ {+}) = \frac {1}{2} (\delta + \mathrm{Ch} (\Lambda ^ {l} \sigma)) \qquad \mathrm{Ch} (\tau^ {-}) = \frac {1}{2} (- \delta + \mathrm{Ch} (\Lambda ^ {l} \sigma)).
\]

\begin{enumerate}
\def\labelenumi{\alph{enumi})}
\setcounter{enumi}{4}
\tightlist
\item
  证明同态 \(u : Z[X_{1}, \ldots, X_{l}; Y] \to R(G)\)（满足 \(u(X_{i}) = [\Lambda^{i} \sigma]\) 与 \(u(Y) = [\tau^{+}]\)）是满射，且其核由 \(Y^{2} - YX_{l} + A\) 生成，其中 \(A \in Z[X_{1}, \ldots, X_{l}]\)。
\end{enumerate}

\begin{enumerate}
\def\labelenumi{\arabic{enumi})}
\setcounter{enumi}{3}
\tightlist
\item
  设 E 是一个复向量空间，\(\mathbf{T}_{q}^{p}(\mathrm{E})\) 是 E 上 \((p,q)\) 型张量组成的空间（《代数》第 III 章，§5，节 6）。用 \(\mathrm{H}_{q}^{p}(\mathrm{E})\) 表示 \(\mathbf{T}_{q}^{p}(\mathrm{E})\) 中由对称张量（即属于 \(\mathbf{T}\mathbf{S}^{p}(\mathrm{E})\otimes\mathbf{T}\mathbf{S}^{q}(\mathrm{E}^{*})\) 在 \(\mathrm{T}_{q}^{p}(\mathrm{E})\) 中的像、且被收缩 \(c_{j}^{i}\)（对 \(i\in\{1,p\}\) 与 \(j\in\{p+1,p+q\}\)）零化的那些张量）组成的子空间（同上）。
\end{enumerate}

\begin{enumerate}
\def\labelenumi{\alph{enumi})}
\item
  证明 \(\mathrm{H}_{q}^{p}(\mathbf{C}^{n})\) 在 \(\mathbf{GL}(n,\mathbf{C})\) 下稳定，从而在 \(\mathbf{SU}(n,\mathbf{C})\) 下稳定；把 \(\mathbf{SU}(n,\mathbf{C})\) 的这个表示记作 \(\tau_{q}^{p}\)。
\item
  证明 \(\tau_q^p\) 是一个不可约表示，其最高权（用 §4 习题 1 的记号）是 \(p\varpi_1 + q\varpi_{n - 1}\)。
\item
  \(\mathbf{SU}(3,\mathbf{C})\) 的每个不可约表示都同构于诸表示 \(\tau_q^p\) 之一。
\end{enumerate}

\(d\) ) 设 \((x_{i})_{1\leq i\leq n}\) 是 \(\mathbf{C}^n\) 的典则基，\((y_{j})_{1\leq j\leq n}\) 是对偶基。证明 \(\mathrm{H}_q^p (\mathbf{C}^n)\) 可与由多项式 \(\mathrm{P}\in \mathbf{C}[x_1,\ldots ,x_n,y_1,\ldots ,y_n]\)（次数为 \(p\)，关于诸 \(x_{i}\)；次数为 \(q\)，关于诸 \(y_{i}\)；且满足 \(\sum_{i = 1}^{n}\frac{\partial^2\mathrm{P}}{\partial x_i\partial y_i} = 0\)）组成的空间等同。

\begin{enumerate}
\def\labelenumi{\arabic{enumi})}
\setcounter{enumi}{4}
\tightlist
\item
  设 \(k\) 是一个特征为零的交换域，\(V\) 是 \(k\) 上的一个向量空间，而 \(\varPsi\) 是 \(V\) 上一个分离的二次型。设 \(\varGamma \in \mathbf{S}^2(V^*)\)（相应地 \(\varGamma^* \in \mathbf{S}^2(V)\)）是伴随于 \(\varPsi\) 的元素（相应地 \(\varPsi\) 的逆型）。用 \(Q\) 表示 \(S(V)\) 上取与 \(\varGamma^*\) 的乘积所给出的自同态，用 \(\varDelta\) 表示 \(S(V)\) 上取与 \(\varGamma\) 的内积所给出的自同态，用 \(h\) 表示 \(S(V)\) 上在 \(S^r(V)\) 上约为乘以 \(-\frac{n}{2} - r\) 的自同态。\(a\) ) 若 \((x_i)_{1\leq i\leq n}\) 是 \(V\) 的一个标准正交基，则 \(Q(P) = \frac{1}{2}\left(\sum x_i^2\right).P\) 与 \(\varDelta(P) = \frac{1}{2}\sum \frac{\partial^2P}{\partial x_i^2}\)（对 \(P \in S(V)\)）。
\end{enumerate}

\begin{enumerate}
\def\labelenumi{\alph{enumi})}
\setcounter{enumi}{1}
\item
  证明公式 \([\Delta, \mathrm{Q}] = -h\)，\([h, \Delta] = 2\Delta\)，\([h, \mathrm{Q}] = -2\mathrm{Q}\)。
\item
  设 \(H_{r}\) 是 \(\mathbf{S}^{r}(\mathrm{V})\) 中被 \(\Delta\) 零化的子空间（「r 次齐次调和多项式」）。由 b) 推出存在一个在 \(\mathbf{O}(\varPsi)\) 下稳定的直和分解
\end{enumerate}

\[
\mathbf {S} ^ {r} (\mathrm{V}) = \mathrm{H} _ {r} \oplus \mathrm{QH} _ {r - 2} \oplus \mathrm{Q} ^ {2} \mathrm{H} _ {r - 4} \oplus \dots .
\]

\(d\) ) 证明表示 \(\mathrm{H}_r\) 是不可约的（参见~第 VIII 章，§13，节 3，(IV)）。

\begin{enumerate}
\def\labelenumi{\alph{enumi})}
\setcounter{enumi}{4}
\tightlist
\item
  取 \(k = \mathbf{C}\)，\(V = \mathbf{C}^n\) 带通常的二次型（\(n \geq 3\)）。我们于是得到不可约表示 \(\tau_r\)（\(\mathbf{SO}(n, \mathbf{R})\) 的）；用 §4 习题 2 的记号，证明 \(\tau_r\) 的最高权是 \(r\varpi_1\)。
\end{enumerate}

\(f\) ) 设 \(\Gamma\) 是由 Killing 型得到的 G 的 Casimir 元。证明公式 \(\Gamma_{\mathbf{S}(\mathrm{V})} = \frac{1}{2n - 4}\left(-4\mathrm{Q}\Delta +\left(\mathrm{H} + \frac{n}{2}\mathrm{I}\right)\left(\mathrm{H} + \left(2 - \frac{n}{2}\right)\mathrm{I}\right)\right)\)。计算 \(\tilde{\Gamma} (\tau_r)\) 并由此推出型 \(Q_{\Gamma}\) 的值（命题 4）。

\begin{enumerate}
\def\labelenumi{\arabic{enumi})}
\setcounter{enumi}{5}
\item
  假定 G 是概单的。证明 G 具有一个忠实不可约表示当且仅当它不同构于 \(\mathbf{Spin}(4k, \mathbf{R})\)（\(k \geq 2\)）。
\item
  设 \(\tau : \mathrm{G} \to \mathbf{SO}(n, \mathbf{R})\) 是 \(\mathrm{G}\) 的一个实酉表示；用 \(\varphi : \mathbf{Spin}(n, \mathbf{R}) \to \mathbf{SO}(n, \mathbf{R})\) 表示典则二重覆盖。于是称 \(\tau\) 为旋量的，若存在一个态射 \(\tilde{\tau} : \mathrm{G} \to \mathbf{Spin}(n, \mathbf{R})\) 使 \(\varphi \circ \tilde{\tau} = \tau\)。
\end{enumerate}

\begin{enumerate}
\def\labelenumi{\alph{enumi})}
\item
  设 \(\Sigma\) 是 \(\mathrm{P}(\mathrm{T},\tau)\) 的一个满足 \(\Sigma\cup(-\Sigma)=\mathrm{P}(\mathrm{T},\tau)\) 与 \(\Sigma\cap(-\Sigma)=\varnothing\) 的子集；用 \(\omega_{\Sigma}\) 表示 \(\Sigma\) 的诸元素之和。类 \(\varpi\)（\(\omega_{\Sigma}\) 在 \(\mathrm{X}(\mathrm{T})/2\mathrm{X}(\mathrm{T})\) 中的）不依赖于 \(\Sigma\) 的选取。证明 \(\tau\) 是旋量的当且仅当 \(\varpi=0\)。
\item
  证明 \(\rho \in \mathrm{X}(\mathrm{T})\) 当且仅当伴随表示是旋量的。
\end{enumerate}

\begin{enumerate}
\def\labelenumi{\arabic{enumi})}
\setcounter{enumi}{7}
\tightlist
\item
  设 G 是一个李代数为紧的、且只有有限个连通分支的李群。证明 G 在一个有限维向量空间上具有一个忠实线性表示（把 G 写成一个紧群 K 与一个向量群 N 的半直积；选取 K 在一个有限维向量空间 W 上的一个忠实表示，并把 G 表示为 \(W \oplus N\) 的仿射群的一个子群）。
\end{enumerate}

\exercisesection{\S~8}

\begin{enumerate}
\def\labelenumi{\arabic{enumi})}
\tightlist
\item
  设 \(\mathbf{G} = \mathbf{SU}(2, \mathbf{C})\)，\(\sigma\) 是 G 在 \(C^{2}\) 上的恒等表示。
\end{enumerate}

\begin{enumerate}
\def\labelenumi{\alph{enumi})}
\item
  G 的不可约表示是诸表示 \(\tau^{n} = S^{n}\sigma\)（\(n \geq 0\)）。
\item
  设 \(e_{1}, e_{2}\) 是 \(C^{2}\) 的典则基；证明 \(\tau^{n}\) 在基 \((e_{1}^{i}e_{2}^{n-i})_{0\leq i\leq n}\) 下的系数是满足如下条件的函数 \(\tau_{ij}^{n}\)：对 \(g=\begin{pmatrix}\alpha&\beta\\-\bar{\beta}&\bar{\alpha}\end{pmatrix}\in G\)，有 \(\tau_{ij}^{n}(g)=\frac{(-1)^{i}}{j!}\alpha^{i+j-n}\beta^{j-i}\mathrm{P}_{ij}^{n}(|\alpha|^{2})\)，其中 \(\mathrm{P}_{ij}^{n}(t)=\frac{d^{j}}{dt^{j}}\left[t^{n-i}(1-t)^{i}\right]\)（「Jacobi 多项式」）。
\item
  由 \(b)\) 推出，诸函数 \((n + 1)^{1/2}\left(\frac{j!(n - j)!}{i!(n - i)!}\right)^{1/2}\tau_{ij}^n\)（对整数 \(i,j,n\)，满足 \(0\leq i\leq n\)，\(0\leq j\leq n\)）构成 \(\mathrm{L}^2 (\mathrm{G})\) 的一个标准正交基。
\end{enumerate}

\(d)\) 对 \(g = \left( \begin{array}{cc}\alpha & \beta \\ -\bar{\beta} & \bar{\alpha} \end{array} \right)\in \mathrm{G}\)，置 \(\alpha = t^{1 / 2}e^{i(\varphi -\psi) / 2},\beta = (1 - t)^{1 / 2}e^{i(\varphi +\psi) / 2}\)，其中 \(0\leq t\leq 1,0\leq \varphi < 2\pi , - 2\pi \leq \psi < 2\pi\)。证明规范化的 Haar 测度 \(dg\)（\(\mathrm{G}\) 上的）等于 \((8\pi^2)^{-1}dtd\varphi d\psi\)。

\begin{enumerate}
\def\labelenumi{\alph{enumi})}
\setcounter{enumi}{4}
\tightlist
\item
  设 \(a, b\)（属于 \(\frac{1}{2}\mathbf{Z}\)）满足 \(a - b \in \mathbf{Z}\)。由 \(d\) ) 推出，多项式 \(\mathrm{P}_{n/2 - a, n/2 - b}^n(t)\)（\(n\) 为与 \(2a\) 同奇偶的整数，且 \(n \geq \max(2a, 2b)\)）构成 \(\mathrm{L}^2([0,1])\) 关于测度 \(t^{-a - b}(1 - t)^{a - b}dt\) 的一个正交基。
\end{enumerate}

\begin{enumerate}
\def\labelenumi{\arabic{enumi})}
\setcounter{enumi}{1}
\item
  设 \(f\) 是 \(G\) 上的一个 \(C^\infty\) 类复值函数。证明存在两个复值函数 \(g\) 与 \(\varphi\)（\(G\) 上的，\(C^\infty\) 类），使得 \(\varphi\) 是中心的且 \(f = g * \varphi\)。
\item
  设 u 是 G 的一个不可约表示，\(\lambda\) 是它的最高权。对 \(x \in g\)，用 \(\bar{x}\) 表示 \(\overline{C}\) 中在 \(\mathrm{Ad}(G)\) 下与 x 共轭的唯一元素。证明等式 \(\|u(x)\|_{\infty} = |\delta(\lambda)(\bar{x})|\)。
\item
  选取一个分离的正二次型 \(\mathbf{Q}\)（\(\mathfrak{g}\) 上的，在 \(\operatorname{Ad}(\mathbf{G})\) 下不变）。对 \(x \in \mathfrak{t}\)，置 \(\vartheta_0(x) = \sum_{u \in \Gamma(\mathbf{T})} e^{-\mathbf{Q}(x + u)}\)。
\end{enumerate}

\begin{enumerate}
\def\labelenumi{\alph{enumi})}
\item
  证明 \(\vartheta_0\) 是一个 \(C^\infty\) 类函数（在 \(\mathfrak{t}\) 上），且存在函数 \(\vartheta_1\)（\(C^\infty\) 类的，在 \(T\) 上）使 \(\vartheta_1 \circ \exp_T = \vartheta_0\)。
\item
  证明存在唯一的中心函数 \(\vartheta\)（\(\mathbf{G}\) 上的，\(\mathrm{C}^{\infty}\) 类），它在 \(\mathrm{T}\) 上的限制等于 \(\vartheta_{1}\)。对每个极大环面 \(\mathrm{S}\)（\(\mathrm{G}\) 的）与所有 \(x\in \mathrm{L}(\mathrm{S})\)，我们有
\end{enumerate}

\[
\vartheta (\exp x) = \sum_ {u \in \Gamma (S)} e ^ {- Q (x + u)}.
\]

\begin{enumerate}
\def\labelenumi{\alph{enumi})}
\setcounter{enumi}{2}
\tightlist
\item
  设 A 是 t 的一个单房，dx 是 t 上的一个 Haar 测度，h 是 t 上在仿射 Weyl 群 \(W_{a}^{\prime}\) 下不变的局部可积函数（§5，节 2）。证明等式
\end{enumerate}

\[
\int_ {\mathrm{A}} h (x) d x = \frac {1}{w (\mathrm{G})} \int_ {\mathfrak {t}} h (x) e ^ {- \mathrm{Q} (x)} (\vartheta_ {0} (x)) ^ {- 1} d x.
\]

\(d\) ) 对 \(x\in \mathfrak{g}\)，置 \(\xi (x) = \lambda_{\mathfrak{g}}(x)e^{-\mathrm{Q}(x)}(\vartheta (\exp x))^{-1}\)，其中 \(\lambda_{\mathfrak{g}}(x) = \det \frac{e^{\operatorname*{ad}x} - 1}{\operatorname*{ad}x}\)（§6，节 3）。证明 \(\xi\) 是一个 \(\mathrm{C}^\infty\) 类函数（\(\mathfrak{g}\) 上的），且若 \(\mu\) 是 \(\mathfrak{g}\) 上的一个 Haar 测度，则 \(\exp_{\mathrm{G}}\) 下测度 \(\xi \mu\) 的像是 G 上的一个 Haar 测度（利用 \(c\) ) 以及定理 1 的推论 2 与命题 4（§6，节 2 与节 3））。

\begin{enumerate}
\def\labelenumi{\alph{enumi})}
\setcounter{enumi}{4}
\tightlist
\item
  证明公式 \(\vartheta_0(x) = m\sum_{\lambda \in \mathrm{X(T)}}\exp (\delta (\lambda)(x) - \frac{1}{4}\mathbf{Q}'(\delta (\lambda)))\)（对 \(x\in \mathfrak{t}\)），其中 \(\mathbf{Q}'\) 是 \(\mathfrak{t}_{\mathbf{C}}^{*}\) 上的二次型，它是 \(\mathfrak{t}_{\mathbf{C}}\) 上由 \(\mathbf{Q}\) 诱导的二次型的逆型，\(m\) 是一个应当算出的常数（利用 Poisson 公式，参见~《谱理论》（编纂中）第 II 章，§1，节 8）。推出等式 \(\vartheta (t) = m\sum_{\lambda \in \mathrm{X(T)}}e^{-\mathrm{Q}'(\delta (\lambda)) / 4}t^{\lambda}\)（对 \(t\in \mathrm{T}\)）。
\end{enumerate}

\begin{enumerate}
\def\labelenumi{\arabic{enumi})}
\setcounter{enumi}{4}
\tightlist
\item
  \begin{enumerate}
  \def\labelenumii{\alph{enumii})}
  \tightlist
  \item
    设 V 是一个实向量空间，f 是 V 上的一个非零线性形式，H 是 f 的核。~那么，函数 \(\varphi\)（V 上的，\(C^{\infty}\) 类）在 H 上为零当且仅当它可写成 \(f\varphi'\)，其中 \(\varphi'\) 是 V 上的一个 \(C^{\infty}\) 类函数。
  \end{enumerate}
\end{enumerate}

\begin{enumerate}
\def\labelenumi{\alph{enumi})}
\setcounter{enumi}{1}
\item
  设 V 是一个实向量空间，\((f_{i})_{i\in\mathrm{I}}\) 是一个由非零线性形式组成的有限族，使得诸 \(H_{i}=\operatorname{Ker}f_{i}\) 两两不同。那么，函数 \(\varphi\)（V 上的，\(C^{\infty}\) 类）在诸 \(H_{i}\) 的并上为零当且仅当它可写成 \(\psi\prod_{i\in I}f_{i}\)，其中 \(\psi\) 是 V 上的一个 \(C^{\infty}\) 类函数。
\item
  设 T 是一个环面，\((\alpha_{i})_{i\in I}\) 是一个由 T 的异于 1 的特征标组成的有限族，使得诸 \(K_{i} = Ker \alpha_{i}\) 两两不同。那么，函数 \(\varphi\)（T 上的，\(C^{\infty}\) 类）在诸 \(K_{i}\) 的并上为零当且仅当它可写成 \(\psi \cdot \prod_{i\in I}(\alpha_{i}-1)\)，其中 \(\psi\) 是 T 上的一个 \(C^{\infty}\) 类函数（在 T 上局部地论证并利用 b)）。
\end{enumerate}

\(d\) ) 用节 4 的记号，证明映射 \(b_{\infty}:\mathcal{C}^{\infty}(\mathrm{T})^{\mathrm{W}}\to \mathcal{C}^{\infty}(\mathrm{T})^{-\mathrm{W}}\) 是双射。

\begin{enumerate}
\def\labelenumi{\arabic{enumi})}
\setcounter{enumi}{5}
\item
  假定 G 不是交换的。证明 T 上的连续函数 \(\mathrm{J}(\rho)^{1/3}\) 在 W 下是反不变的，但不属于映射 \(b_{c}: \mathcal{C}(\mathrm{T})^{\mathrm{W}} \to \mathcal{C}(\mathrm{T})^{-\mathrm{W}}\) 的像。
\end{enumerate}

\exercisesection{\S~9}

\begin{enumerate}
\def\labelenumi{\arabic{enumi})}
\item
  设 A 是 R 的由 0 与诸实数 \(1/n\)（\(n\) 为 \(\geq 1\) 的整数）组成的紧子集。证明当 \(\mathcal{C}^r(\mathbf{R};\mathbf{R})\) 赋予 A 上的一致 \(\mathrm{C}^r\) 收敛拓扑时，在 A 的邻域中为嵌入的态射组成的集合在 \(\mathcal{C}^r(\mathbf{R};\mathbf{R})\) 中不是开的（考察一个函数序列 \((f_n)_{n\geq 1}\)，使得 \(f_n(x) = x\)（对 \(x\leq \frac{1}{n + 1}\)），\(f_n(x) = x - \frac{1}{n}\)（对 \(x\geq \frac{1}{n}\)））。\\
\item
  设 X 是一个 \(\mathrm{C}^r\) 类（\(1\leq r\leq \infty\)）的分离流形，它无穷可数且是维数 \(n\) 的纯流形。
\end{enumerate}

\begin{enumerate}
\def\labelenumi{\alph{enumi})}
\item
  假定存在 X 到一个有限维向量空间 V 中的嵌入 \(\varphi\)。证明存在 X 在 \(\mathbf{R}^{2n + 1}\) 中的一个嵌入（若 \(\dim \mathrm{V} > 2n + 1\)，证明存在一点 \(p\)（V 中的），使得对所有 \(x \in \mathrm{X}\)，连接 \(p\) 与 \(\varphi(x)\) 的直线与 \(\varphi(\mathrm{V})\) 只交于点 \(\varphi(x)\)，且在该点横截相交；由此推出存在 X 在一个维数等于 \(\dim \mathrm{V} - 1\) 的空间中的嵌入）。
\item
  证明存在 X 在 \(R^{2n+1}\) 中的一个嵌入。（设 O 是 X 的开子集组成的集合，U 是 O 中由这样的开集 U 组成的子集：存在一个在 U 上的限制为嵌入的态射 \(\varphi: X \to R^{2n+1}\)；利用 a)，证明 U 是 O 的一个拟全子集（《一般拓扑学》第 IX 章，§4，习题 27），从而等于 O。）
\item
  证明存在 X 在 \(\mathbf{R}^{2n + 1}\) 中的一个真嵌入。（借助 X 上的一个真函数构造 X 在 \(\mathbf{R}^{2n + 2}\) 中的一个真嵌入。）
\end{enumerate}

¶ 3) 设 G 是一个李群，H 是 G 的一个紧子群。假定 G 具有一个忠实的（有限维）线性表示。

\begin{enumerate}
\def\labelenumi{\alph{enumi})}
\item
  设 \(\Theta_{\mathrm{H}}(\mathrm{G})\) 是 \(\mathcal{C}(\mathrm{H};\mathbf{R})\) 中由 G 上的（连续）代表函数在 H 上的限制组成的子代数。证明 \(\Theta_{\mathrm{H}}(\mathrm{G})\) 在 \(\mathcal{C}(\mathrm{H};\mathbf{R})\) 中对一致收敛拓扑是稠密的。
\item
  设 \(f \in \Theta_{\mathrm{H}}(\mathrm{G})\)。证明存在 G 在一个有限维实向量空间上的表示 \(\sigma\)，使得表示 \(\sigma|H\) 的某个系数与 f 不正交（《谱理论》（编纂中））。
\item
  设 \(\rho:H\to\mathbf{GL}(V)\) 是 H 在一个有限维实向量空间上的表示。证明存在一个有限维实向量空间 W、一个表示 \(\sigma:G\to\mathbf{GL}(W)\) 与一个单同态 \(u:V\to W\)，使得 \(u(\rho(h)v)=\sigma(h)u(v)\)（对 \(h\in H\)，\(v\in V\)）。
\end{enumerate}

（化归到 \(\rho\) 不可约的情形，并利用 b)。）

\begin{enumerate}
\def\labelenumi{\arabic{enumi})}
\setcounter{enumi}{3}
\item
  设 G 是一个李群，H 是 G 的一个紧子群，\(\rho\) 是 H 在一个实希尔伯特空间 V 上的一个酉表示。证明存在一个实希尔伯特空间 W、G 在 W 上的一个酉表示 \(\sigma\) 以及一个具有闭像的等距单射 \(u: V \to W\)，使得 \(u(\rho(h)v) = \sigma(h)u(v)\)（对 \(h \in H\)，\(v \in V\)）（方法与习题 3 相同）。
\item
  设 G 是一个李群，H 是 G 的一个紧子群。假定存在 G 在一个有限维实向量空间上的一个忠实线性表示 \(\rho: G \to \mathbf{GL}(V)\)。
\end{enumerate}

\begin{enumerate}
\def\labelenumi{\alph{enumi})}
\item
  证明存在表示 \(\sigma\)（\(\mathbf{GL}(\mathrm{V})\) 在一个有限维实向量空间 W 上的）、V 上一个分离的正二次型 q 以及 W 中的一个向量 w，使得 \(\rho(\mathrm{H}) = \mathbf{O}(q) \cap \mathrm{F}_{w}\)，其中 \(F_{w}\) 是 w 在 \(\mathbf{GL}(\mathrm{V})\) 中的稳定子群（选取一个在 H 下不变的二次型 q 以及 \(\mathbf{O}(q)\) 的一个使 \(\rho(\mathrm{H})\) 为一点的稳定子群的表示（节 2 的推论 2），然后应用习题 3 c)）。
\item
  由 a) 推出，存在一个有限维实向量空间 E、G 在 E 上的一个表示以及一个以 H 为稳定子群的向量 \(e \in E\)（取 \(E = W \oplus Q\)，其中 Q 是 V 上二次型组成的空间，且 \(e = (w, q)\)）。
\end{enumerate}

\begin{enumerate}
\def\labelenumi{\arabic{enumi})}
\setcounter{enumi}{5}
\item
  设 G 是一个李群，H 是 G 的一个紧子群。证明存在 G 在一个实希尔伯特空间 E 上的一个酉表示以及一个以 H 为稳定子群的向量 \(e \in E\)（取 \(\mathrm{E} = \mathrm{L}^{2}(\mathrm{G})\)）。
\item
  设 G 是一个李群，H 是 G 的一个紧子群，\(\rho\) 是 H 在一个实希尔伯特空间 W 上的一个酉表示。用 X 表示流形 \(G \times ^{H}W\)（节 3）。
\end{enumerate}

\begin{enumerate}
\def\labelenumi{\alph{enumi})}
\item
  证明存在一个希尔伯特空间 V、G 在 V 上的一个酉表示 \(\sigma\) 以及一个（解析）嵌入 \(\varphi: X \to V\)，使得 \(\varphi(gx) = \sigma(g)\varphi(x)\)（对 \(g \in G\)，\(x \in X\)）（利用习题 4 与习题 6）。
\item
  证明若 W 是有限维的且 G 具有一个忠实的有限维线性表示，则 V 可取为有限维的（利用习题 3 与习题 5）。
\end{enumerate}

¶ 8) 设 X 是一个 \(C^{r}\) 类（\(1 \leq r \leq \infty\)）的仿紧流形，G 是一个真地作用在 X 上的李群，且作用律 \((g, x) \mapsto gx\) 是 \(C^{r}\) 类的。

\begin{enumerate}
\def\labelenumi{\alph{enumi})}
\item
  证明存在 G 在一个希尔伯特空间 V 上的一个酉表示 \(\rho\) 以及一个嵌入 \(\varphi\)（\(C^{r}\) 类的，X 到 V 中的），使得 \(\varphi(gx) = \sigma(g)\varphi(x)\)（对 \(g \in G, x \in X\)）。（利用命题 6、习题 7，并按命题 4 的证明方式论证。）
\item
  再作如下补充假定：
\end{enumerate}

\begin{enumerate}
\def\labelenumi{(\roman{enumi})}
\item
  G 具有一个有限维线性表示；
\item
  \(\mathrm{X}\) 无穷可数且维数有界；
\item
  X 在 G 的作用下只有有限多个轨道类型。
\end{enumerate}

证明存在 X 的一个由在 G 下稳定的开集 \((\mathrm{U}_{i})_{i\in\mathrm{I}}\) 组成的有限覆盖、G 的紧子群 \((\mathrm{H}_{i})_{i\in\mathrm{I}}\)，以及对每个 \(i\in I\)，一个闭子流形 \(S_{i}\)（\(U_{i}\) 的，在 \(H_{i}\) 下稳定），使得映射 \((g,s)\mapsto gs\) 通过过渡到商空间诱导出从 \(G\times^{H_{i}}S_{i}\) 到 \(U_{i}\) 的一个同构（证明 X/G 中其原像在 X 内容许这样一个覆盖的开子集组成 X/G 的开子集集合的一个拟全子集，参见~《一般拓扑学》第 IX 章，§4，习题 27）。

\begin{enumerate}
\def\labelenumi{\alph{enumi})}
\setcounter{enumi}{2}
\tightlist
\item
  在 b) 的假定下，证明存在 G 在一个有限维实向量空间 V 上的一个表示以及 X 在 V 中的一个与 G 的作用相容的 \(C^{r}\) 类嵌入（在 G 的子群组成的集合上用归纳法论证，利用 b)、习题 2 与习题 7）。
\end{enumerate}

\begin{enumerate}
\def\labelenumi{\arabic{enumi})}
\setcounter{enumi}{8}
\tightlist
\item
  设 \(G\) 是一个真地作用在流形 \(X\) 上的李群。作下列两个假定之一：
\end{enumerate}

\begin{enumerate}
\def\labelenumi{(\roman{enumi})}
\tightlist
\item
  X/G 是紧的，且 X 无穷可数、维数有界；\\
\item
  X 是一个 G 在其上线性作用的有限维向量空间。
\end{enumerate}

证明 X 的元素的轨道类型组成的集合是有限的（同时对两种情形按 dim X 用归纳法处理，并利用命题 6）。

\[
\mathbf {G L} (3, \mathbf {R})
\]

\(\left(\begin{array}{ccc}1 & a & b \\ 0 & 1 & 0 \\ 0 & 0 & 1\end{array}\right), a, b \in \mathbf{R}.\) 证明，对 \(G\) 在

\(\mathbf{R}^3\) 上的恒等表示，轨道类型的个数是无限的。

\end{enumerate}
